\documentclass[11pt]{amsart}

% Encoding and fonts
\usepackage[T1]{fontenc}
\usepackage[utf8]{inputenc}
\usepackage{lmodern}

% Math packages
\usepackage{amsmath}
\usepackage{amssymb}
\usepackage{amsfonts}
\usepackage{amsthm}
\usepackage{mathtools}

% Layout and utilities
\usepackage{geometry}
\usepackage{enumitem}
\usepackage{array}
\usepackage[hidelinks]{hyperref}

\geometry{margin=1in}

\numberwithin{equation}{section}
\setcounter{tocdepth}{2}
\setcounter{secnumdepth}{3}

% ---------------------------------------------------------
% Theorem-like environments
% ---------------------------------------------------------

\theoremstyle{plain}
\newtheorem{theorem}{Theorem}[section]
\newtheorem{proposition}[theorem]{Proposition}
\newtheorem{lemma}[theorem]{Lemma}
\newtheorem{corollary}[theorem]{Corollary}
\newtheorem{conjecture}[theorem]{Conjecture}

\theoremstyle{definition}
\newtheorem{definition}[theorem]{Definition}
\newtheorem{experiment}[theorem]{Experiment}
\newtheorem{heuristic}[theorem]{Heuristic}
\newtheorem{question}[theorem]{Question}
\newtheorem{objective}[theorem]{Objective}

\theoremstyle{remark}
\newtheorem{remark}[theorem]{Remark}
\newtheorem{warning}[theorem]{Warning}

\DeclareMathOperator{\tr}{tr}
\DeclareMathOperator{\rank}{rank}
\DeclareMathOperator{\Var}{Var}
\DeclareMathOperator{\Cov}{Cov}
\DeclareMathOperator{\rad}{rad}
\DeclareMathOperator*{\esssup}{ess\,sup}

\title{
Orientation-Bit Geometry of the M\"obius Path:
Reflection Fibres, Divisor--Shift Boundary Cocycles,
and an Observability Route to Two-Point Cancellation
}

\author{}
\date{Version 6}

\begin{document}

\begin{abstract}
This notebook develops an operator framework for the two-point
correlation $C_N(h)=\sum_{n\le N-h}\mu(n)\mu(n+h)$ built on the
orientation bits of the M\"obius path and on the divisor operator.

Version 3 adds an analytic layer to the combinatorial one. The central
new object is the difference $(U_h-I)D$ rather than the commutator or
the conjugation: each divisor column collapses to the difference of two
residue classes, vanishes identically when the divisor divides $h$, and
has interval sums bounded by one. A purely algebraic Haar basis over an
arbitrary binary tree then annihilates the volumetric CRT term exactly,
leaving only boundary cocycles which satisfy a parent--children additivity
law. No doubling, no comparable lengths and no regularity of the
squarefree mask are required at any point.

Four further results are proved. An exact observability criterion:
$|x^{\mathsf T}Ax|\le\kappa\|Lx\|^2$ holds for some finite $\kappa$ if
and only if $\ker L\subseteq\ker A$, with the optimal constant and the
exact operator-norm dual given in closed form. A no-go: uniform
observability of the squarefree shift by a sublinear set of divisor rows
is impossible, because above $N/2$ each column is seen only by its own
row and shifted squarefree pairs have positive density. An extension of
the blindness theorem to kernels that are constant on each half of an
excursion, under which the cross-excursion interaction factors through
the single channel $s_jH_j$ and the shift kernel is shown not to belong
to the extended class. And a geometric proposition: the hypothesis
$u_j\ge m_j/2$ needed for the central divisor matrix to be triangular
with nonvanishing diagonal fails for at most $\lfloor\log_2 N\rfloor$
excursions, so that the M\"obius line is isolated up to a kernel of
logarithmic dimension.

The results of Versions 1 and 2 are retained: the reflection fibre of
the quadratic profile, the exactly isotropic null, the blindness theorem
for kernels depending on $\max(a,b)$, and the criterion separating
tautological from informative sine-kernel content. Negative results,
including the collapse of every multiplicative diagonal weight onto
$1/\zeta(1+s)$, are recorded in full. No numerical computation is
performed.

Version 4 audits the target objective of Version 3 and finds it
self-defeating as stated. An exact identity expresses the $L^2$ defect
of any leaf projection in the squarefree space as an explicit
combination of within-leaf correlations of $\mu$ along consecutive
squarefree integers. The capture condition of the adaptive boundary
objective therefore demands maximally positive correlations on every
nontrivial leaf, the opposite of the cancellation it is meant to
produce, while long leaves are excluded unconditionally by
short-interval cancellation. The objective is marked superseded. The
engine common to both objectives, the annihilation $D\mu=e_1$, survives
in the central-row objective, which becomes the live target. The
conditioning of the two-point expansion on central data and the
redundancy of quarter and octave evaluations are also recorded.

Version 5 incorporates an external audit, verified line by line. It
corrects the pushed form to $\widehat G=F^{\mathsf T}GF$ with incremental
factors, fixes a single orientation basis, states the conditioning on
$M(N)$ that freezes the terminal block, corrects the parity statement for
single reflections, and withdraws two overstatements (a semicircular
density inferred from an edge bound, and a general dichotomy for
non-multiplicative weights). It withdraws the orthogonal-to-unitary
crossover conjecture as stated: a Hermitian deformation supported on a
forest, or obtained by unitary transport, is isospectral, so a crossover
needs flux around cycles. It adds four exact results: the energy kernel
is Montgomery's weight localised at $\alpha=0$ with tail
$\Pr(|\alpha|>A)=T^{-2A}$; the Hermitian collision constant differs from
the M\"obius energy constant by an exact fourth-order term; the form
$\mathcal A_0'$ is the discrete L\'evy area of the planar path
$(M_\Omega,L)$, neither a product nor a fixed-shift correlation; and the
central divisor rows cut the reflection fibre from $2^Z$ orientations to
at most $2N$, with at most $\lfloor\log_2N\rfloor+1$ further divisor
rows selecting the M\"obius orientation uniquely.

Version 6 executes the uncompressed, irregular route analytically and
closes it as a family. The interior divisor--shift operator
$B_h^\circ=R_{N-h}(U_h-I)D$ has maximal rank $N-h$: the vanishing of the
columns $d\mid h$ removes only $\tau(h)$ columns. Each nonzero column is
the exact discrete derivative of a periodic block of $h\bmod d$
consecutive residues, which yields an optimal variation inequality on
arbitrary intervals, but the Haar energy of each column is independent
of the tree. The quantitative no-go is sharpened to
$|\Lambda|\ge\frac{\mathfrak S_h}{4}N+o_h(N)$ divisor rows for any
uniform observability with error below $1/2$. The shift does not descend
to valuation-indexed local factors. The maximal class of kernels blind
to zero-sum blocks is exactly the additively separable class, the shift
lies at macroscopic distance from it, and no separable approximation of
sublinear rank is close in operator norm. What remains open inside this
family is a collective signed cancellation of CRT boundary errors with
M\"obius-linked coefficients, which already requires control of the
parity of large primes.
\end{abstract}

\maketitle

\tableofcontents

\bigskip

\begin{remark}
\label{rem:exploratory_nature}
This notebook is exploratory. Rigorous statements appear inside theorem,
proposition, lemma or corollary environments with proofs. Results
established in separate work are collected in
Section~\ref{sec:standing_results} and stated without proof. Everything
else is labelled heuristic, conjectural or objective. No numerical
computation is performed; the empirical values in
Remark~\ref{rem:empirical_inputs} are inputs from separate work.
\end{remark}

\begin{warning}
\label{warn:bibliographic_metadata}
Bibliographic metadata in Section~\ref{sec:bibliography} has not been
verified against primary sources. DOIs are omitted where not confirmed.
\end{warning}

\section{Introduction}
\label{sec:introduction}

\subsection{Objects}
\label{subsec:objects}

Let $\mu$ be the M\"obius function, $q(n)=\mu(n)^2$,
\begin{equation}
M(x)=\sum_{n\le x}\mu(n),
\qquad
C_N(h)=\sum_{n\le N-h}\mu(n)\mu(n+h).
\label{eq:mertens_and_correlation}
\end{equation}
Let $D$ be the divisor operator $(Dx)(r)=\sum_{d\mid r}x(d)$ and $U_h$
the truncated shift $(U_hx)(n)=x(n+h)$, with symmetric part
$A_h=(U_h+U_h^{*})/2$, so that
\begin{equation}
C_N(h)=\mu^{\mathsf T}A_h\mu .
\label{eq:correlation_as_form}
\end{equation}

\subsection{What is proved in this notebook}
\label{subsec:what_is_proved}

\begin{enumerate}[label=(\roman*)]
\item \emph{Reflection fibre} (Prop.~\ref{prop:reflection_fibre}): the
fibre of the quadratic profile is exactly an orbit of
$\mathcal R_N\cong(\mathbb Z/2)^{Z(N)}$.
\item \emph{Exact isotropy} (Prop.~\ref{prop:exact_isotropy}):
$\mathbb E_\nu[ss^{\mathsf T}]=I_Z$ with no correction term, and closed-form
variance.
\item \emph{Blindness} (Thm.~\ref{thm:cross_excursion_vanishing}): every
cross-excursion contribution to a $\max(a,b)$-kernel vanishes
identically.
\item \emph{Extended blindness}
(Thm.~\ref{thm:half_constant_blindness}): for kernels constant on each
half of an excursion the cross-excursion interaction factors through the
rank-one channel $s_jH_j$; the shift kernel is shown not to lie in this
class.
\item \emph{Divisor--shift boundary lemma} (Lem.~\ref{lem:boundary}):
$((U_h-I)D)_{n,d}=\mathbf 1[d\mid n+h]-\mathbf 1[d\mid n]$, identically
zero for $d\mid h$, supported on two residue classes otherwise, with
interval sums bounded by $1$.
\item \emph{Exact CRT pullback} (Lem.~\ref{lem:crt_pullback}):
$(D^{\mathsf T}U_hD)_{d,e}$ vanishes unless $(d,e)\mid h$ and otherwise
equals $(N-h)/[d,e]$ up to an error of modulus at most $1$.
\item \emph{Algebraic Haar and exact bulk vanishing}
(Lem.~\ref{lem:haar_basis}, Thm.~\ref{thm:haar_crt}): on an arbitrary
binary tree the volumetric CRT term is annihilated exactly and only a
boundary cocycle survives, additive over children.
\item \emph{Observability criterion} (Thm.~\ref{thm:observability}):
$|x^{\mathsf T}Ax|\le\kappa\|Lx\|^2$ for some $\kappa<\infty$ iff
$\ker L\subseteq\ker A$; optimal constant and operator-norm dual in
closed form.
\item \emph{Sublinear no-go} (Thm.~\ref{thm:sublinear_nogo}): no
nontrivial observability of the squarefree shift by $o(N)$ divisor rows.
\item \emph{Logarithmic exceptional set}
(Prop.~\ref{prop:log_exceptional}): the geometric hypothesis making the
central divisor matrix triangular with nonvanishing diagonal fails for at
most $\lfloor\log_2 N\rfloor$ excursions.
\item \emph{Central dichotomy} (Lem.~\ref{lem:central_dichotomy}): every
linear reflection-odd functional of an excursion either vanishes on the
whole fibre or reads the orientation bit directly.
\item \emph{Leaf-energy identity and squeeze}
(Prop.~\ref{prop:leaf_energy}, Cor.~\ref{cor:squeeze}): the $L^2$
defect of a leaf projection equals
$\sum_I\bigl(k_I-1-\tfrac{2}{k_I}\sum_{i<i'\in I}\mu(a_i)\mu(a_{i'})\bigr)$,
so $L^2$ capture of $\mu$ by a subspace independent of its signs forces
either essentially the whole squarefree space or maximal positive
within-leaf correlation.
\item \emph{Gauge obstruction} (Prop.~\ref{prop:gauge}): Hermitian
deformations on forests, and unitary transports, are isospectral.
\item \emph{Collision versus energy constant}
(Prop.~\ref{prop:cherm_energy}): exact fourth-order difference.
\item \emph{L\'evy-area form} (Prop.~\ref{prop:levy_area}):
$\mathcal A_0'(N)$ is the discrete signed area of $(M_\Omega,L)$.
\item \emph{Compressed selection} (Cor.~\ref{cor:selection_count},
Cor.~\ref{cor:selection_completion}): at most $2N$ admissible
orientations, and unique selection after $\lfloor\log_2N\rfloor+1$
further divisor rows.
\item \emph{Interior operator} (Thm.~\ref{thm:interior_rank},
Prop.~\ref{prop:residue_boundary}, Prop.~\ref{prop:tree_independent}):
$\rank B_h^\circ=N-h$; columns are discrete derivatives of periodic
residue blocks; Haar energy is independent of the tree.
\item \emph{Blind class} (Thm.~\ref{thm:blind_class}): kernels blind to
zero-sum blocks are exactly $\alpha(a)+\beta(b)$; the shift is at
macroscopic distance.
\end{enumerate}

\subsection{Inputs and the methodological point}
\label{subsec:inputs}

Results established in separate work are listed with status in
Section~\ref{sec:standing_results}; the two used actively are the tensor
factorisation \ref{item:S4} and the permutation-null mean \ref{item:S10}.

The methodological criterion of Version 2 is retained: a statistic
carries sine-kernel content non-tautologically if and only if its value
under the null differs from its value under the sine kernel
(Remark~\ref{rem:nontautological_criterion}). Version 3 adds the
corresponding criterion for weights: every diagonal weight with a
multiplicative symbol collapses onto $1/\zeta(1+s)$
(Thm.~\ref{thm:multiplicative_collapse}), so informative weights must
break the Euler factorisation.

\section{Standing results used as inputs}
\label{sec:standing_results}

Established in separate work by the author, stated without proof, no
novelty claimed here.

\begin{remark}
\label{rem:standing_results_list}
\begin{enumerate}[label=(S\arabic*)]
\item \label{item:S1}
$c\sqrt N\le\lambda_{\max}(D_N)\le C_\varepsilon N^{1/2+\varepsilon}$ with
Cauchy interlacing corollaries. \emph{Context:
Rem.~\ref{rem:negative_result_density_of_states}.}
\item \label{item:S2}
Trace-norm convergence of $G_N/N^2$ to a compact limit operator.
\emph{Context.}
\item \label{item:S3}
Exact Chinese-remainder formula for the entries of $Q_N$.
\emph{Context; source of \ref{item:S4}.}
\item \label{item:S4}
The tensor factorisation
\begin{equation}
\widehat K(r,q)=\prod_p L_p\bigl(v_p(r),v_p(q)\bigr)
\label{eq:tensor_factorisation_standing}
\end{equation}
over all of $\mathbb N\times\mathbb N$, explicit tridiagonal local
factors, sign corrected. \emph{Used actively as (G2).}
\item \label{item:S5}
$C_{\mathrm{herm}}(\rho)=2(\gamma^2-\beta^2)/(|\rho|^4|\zeta'(\rho)|^2)>0$
for every nontrivial zero, without assuming RH. \emph{Context; see
Rem.~\ref{rem:extremal_abscissa}.}
\item \label{item:S6}
$z_0''(n)=-\tfrac{r^2}{2}\Var_{p\mid n}(\log p)$ for squarefree $n$ with
$r$ prime factors, and
$K(s)=\tfrac{\zeta(s)}{2}(S_0S_2-S_1^2)$ with $S_2=(\zeta'/\zeta)'$.
\emph{Context.}
\item \label{item:S7}
Closed form of $z_\vartheta(n)$ as a finite cosine sum over ordered
partitions via the cocycle $\Phi(\pi)$, with
$A_r=3(6r-11)/8$, $B_r=r(6r-11)/8$ exact for all $r$. \emph{Context;
source of the twisted action in Conj.~\ref{conj:twisted_orientation_bits}.}
\item \label{item:S8}
Blindness to $\Theta$ of the whole tower of \emph{fixed} transversal
jets, with complete proof. \emph{Context; the restriction to fixed jets
leaves the scaling window open.}
\item \label{item:S9}
Unconditional equivalence between RH and tail stability of a two-channel
positive logarithmic deformation of the ordinal-prime transform.
\emph{Context.}
\item \label{item:S10}
$\mathbb E_\pi[E_N]/N=1/\pi^2+M(N)^2/(3N)+o(1)$ with
$E_N=\tfrac1N\sum_{n\le N}M(n)^2$. \emph{Used actively.}
\end{enumerate}
\end{remark}

\begin{remark}
\label{rem:extremal_abscissa}
At the extremal abscissa $s=2\Theta$ all contributions come from zeros
with $\beta=\Theta$ and carry the positive coefficient of \ref{item:S5};
a sum of positive terms does not cancel, so cancellation between distinct
pairs cannot occur there. What remains open is the case of a
non-attained supremum, requiring a limiting formulation.
\end{remark}

\section{Three spaces and the maps between them}
\label{sec:three_spaces}

\begin{definition}
\label{def:three_spaces}
Let $A_N=\{a_1<\dots<a_m\le N: q(a_j)=1\}$, $m=|A_N|$. Set
\begin{align}
H_N&=\mathbb R^N &&\text{(full space; }D,\ U_h\text{)},
\label{eq:space_full}\\
H_{\mathrm{sf}}&=\mathbb R^m &&\text{(squarefree space)},
\label{eq:space_sf}\\
H_{\mathrm{exc}}&=\mathbb R^{Z} &&\text{(excursion space)} .
\label{eq:space_exc}
\end{align}
Let $Je_j=e_{a_j}$, so $J^{*}J=I$ and $P=JJ^{*}=\mathrm{diag}(q)$, and
let $F$ be the canonical incremental matrix
\begin{equation}
F_{n,j}=\Delta Q_j(n)\,\mathbf 1_{E_j}(n),
\label{eq:canonical_F}
\end{equation}
with $Q_j=|M|$ on $E_j$ (Definition~\ref{def:excursion_signs}). Throughout,
$s\in\{\pm1\}^Z$ denotes excursion signs in this canonical basis and
$s_\mu$ the signs of the M\"obius path, so that $\mu=Fs_\mu+x_{\mathrm{term}}$.
\end{definition}

\begin{remark}
\label{rem:information_ledger}
The three spaces retain different information.
$H_N$ retains position, triangularity, divisibility and boundary data.
$H_{\mathrm{sf}}$ retains signs and the integer labels $a_j$; if the
labels are discarded, divisibility is lost. $H_{\mathrm{exc}}$ retains
only block orientations: all intrablock variation, gcd classes,
valuations and individual positions are destroyed.
Since $P\mu=\mu$, the inclusion $J$ loses nothing for the vector $\mu$
itself, but it does lose operator information about perturbations off the
mask.
\end{remark}

\begin{proposition}
\label{prop:fibre_expansion}
With $U_h^{\mathrm{sf}}=J^{*}U_hJ$, $A_h=(U_h+U_h^{*})/2$ and
$x=Fs+x_{\mathrm{term}}$ on the fibre,
\begin{equation}
C_N(h)=s^{\mathsf T}(F^{*}A_hF)s
+2s^{\mathsf T}F^{*}A_hx_{\mathrm{term}}
+x_{\mathrm{term}}^{\mathsf T}A_hx_{\mathrm{term}} .
\label{eq:fibre_expansion}
\end{equation}
\end{proposition}

\begin{proof}
Substitute $x=Fs+x_{\mathrm{term}}$ into
\eqref{eq:correlation_as_form} and expand the symmetric form.
\end{proof}

\begin{remark}
\label{rem:exc_is_subordinate}
Equation~\eqref{eq:fibre_expansion} is exact, but by
Cor.~\ref{cor:profile_blindness} the excursion space is a severe
compression and must be subordinate: it is used in
Section~\ref{sec:centers} only for the reflection-odd component and as
one possible family of stopping blocks. The primary space is $H_N$.
\end{remark}

\section{The three-constraint model}
\label{sec:three_conditionings}

\begin{definition}
\label{def:mask}
The orientation bits are
$\varepsilon=(\varepsilon_n)_{n\in A_N}\in\{\pm1\}^m$,
$\varepsilon_n=\mu(n)$. With $m_\pm$ the number of $n\in A_N$ with
$\mu(n)=\pm1$ and $\bar\varepsilon=M(N)/m$, the \emph{fixed-count sign
bridge} $\pi$ is uniform on the vectors of $\{\pm1\}^m$ with $m_+$
entries $+1$.
\end{definition}

\begin{proposition}
\label{prop:bridge_covariance}
Under $\pi$,
\begin{equation}
\mathbb E_\pi[\varepsilon\varepsilon^{\mathsf T}]=(1-\beta)I+\beta J_{m},
\qquad \beta=\frac{m\bar\varepsilon^{\,2}-1}{m-1},
\label{eq:bridge_covariance}
\end{equation}
$J_m$ the all-ones matrix. Consequently, for an orthogonal projection
$P'$ of rank $d$,
\begin{equation}
\mathbb E_\pi[\langle\varepsilon,P'\varepsilon\rangle]
=(1-\beta)d+\beta\|P'\mathbf 1\|^2=d+O(1)
\quad\text{if }M(N)=o(\sqrt m).
\label{eq:bridge_projection_mean}
\end{equation}
\end{proposition}

\begin{proof}
Diagonal entries are $1$. For $n\ne n'$, counting ordered pairs of
distinct positions gives
$\bigl((m_+-m_-)^2-m\bigr)/(m(m-1))$, and $m_+-m_-=m\bar\varepsilon$.
Trace against $P'$ using $\tr P'=d$, $\tr(P'J_m)=\|P'\mathbf 1\|^2$;
for the last claim $|\beta|\le 1/(m-1)$ once $M(N)^2\le m$.
\end{proof}

\begin{remark}
\label{rem:N1_usage}
The factorisation of $\Cov_\pi(M(n)^2,M(n')^2)$ through the three-block
decomposition $(n,n'-n,N-n')$, with its $O(m)$ variance evaluation, is
recorded separately. Only its structural consequence is used: the
covariance kernel is separable of bounded rank. By
Prop.~\ref{prop:exact_isotropy} the third conditioning replaces this by
an exact closed form.
\end{remark}

\section{The reflection fibre}
\label{sec:reflection_fibre}

\begin{definition}
\label{def:excursions}
Let $0=z_0<z_1<\dots<z_Z$ enumerate $\{n\le N:M(n)=0\}\cup\{0\}$, and
\begin{equation}
E_j=\{n:z_{j-1}<n\le z_j\},\qquad j=1,\dots,Z,
\label{eq:excursion_blocks}
\end{equation}
with terminal block $E_{Z+1}=\{n:z_Z<n\le N\}$, empty iff $M(N)=0$.
Write $u_j=z_{j-1}$, $v_j=z_j$. The reflection group
$\mathcal R_N\cong(\mathbb Z/2)^Z$ acts by
$(\sigma\cdot\varepsilon)_n=\sigma_{j(n)}\varepsilon_n$, the terminal
block held fixed. Holding it fixed is a conditioning on $M(N)$, which the
three-constraint model imposes through the fixed sign count.
\end{definition}

\begin{proposition}
\label{prop:reflection_fibre}
$M_{\sigma\cdot\varepsilon}(n)=\sigma_{j(n)}M_\varepsilon(n)$, so the
quadratic profile $\Pi_N=(M(n)^2)_{n\le N}$ is
$\mathcal R_N$-invariant; conversely two sign vectors on the same mask
with the same profile and the same value of $M(N)$ differ by an element
of $\mathcal R_N$. The fibre of $\Pi_N$ conditioned on $M(N)$ is exactly
the orbit. Without that conditioning, if $M(N)\ne0$ the terminal block
can also be reflected and the fibre of $\Pi_N$ is an orbit of
$(\mathbb Z/2)^{Z+1}$.
\end{proposition}

\begin{proof}
For $n\in E_j$, $M_\varepsilon(z_{j-1})=0$ so
$M_\varepsilon(n)=\sum_{z_{j-1}<k\le n}\varepsilon_k$ lies inside $E_j$
and $\sigma$ multiplies it by $\sigma_j$. Conversely equal profiles give
the same zero set, hence the same blocks, and $|M|$ pointwise; inside a
block $M$ has constant sign, and the ratio of the two signs defines
$\sigma_j$; increments determine the path.
\end{proof}

\begin{definition}
\label{def:excursion_signs}
The excursion signs $s\in\{\pm1\}^Z$ record the sign of $M$ in the
interior of each $E_j$. The reflection null $\nu$ is uniform on
$\{\pm1\}^Z$. Write $Q_j(n)=|M(n)|$ for $n\in E_j$, so that
\begin{equation}
M_s(n)=s_jQ_j(n),\qquad \mu_s(n)=s_j\,\Delta Q_j(n),\qquad n\in E_j .
\label{eq:canonical_positive_path}
\end{equation}
\end{definition}

\begin{proposition}
\label{prop:exact_isotropy}
Under $\nu$, $\mathbb E_\nu[ss^{\mathsf T}]=I_Z$ exactly, hence for
symmetric $A$
\begin{equation}
\mathbb E_\nu[\langle s,As\rangle]=\tr A,
\qquad
\Var_\nu[\langle s,As\rangle]=2\sum_{j\ne k}A_{jk}^2,
\label{eq:exact_mean_variance}
\end{equation}
and for a projection $\widehat P$ of rank $d$,
$\mathbb E_\nu=d$ and
$\Var_\nu=2(d-\sum_j\widehat P_{jj}^2)$.
\end{proposition}

\begin{proof}
The $s_j$ are independent Rademacher. Write
$\langle s,As\rangle=\tr A+\sum_{j\ne k}A_{jk}s_js_k$ and use
$\mathbb E_\nu[s_js_ks_{j'}s_{k'}]=\mathbf 1[\{j,k\}=\{j',k'\}]$;
for a projection $\sum_{j,k}\widehat P_{jk}^2=\tr\widehat P=d$.
\end{proof}

\begin{remark}
\label{rem:hierarchy_of_nulls}
Conditioning on the mask removes the $1/\pi^2$ of \ref{item:S10};
conditioning on the count removes the drift $M(N)^2/(3N)$; conditioning
in addition on the profile leaves exactly $\mathcal R_N$, whose null is
isotropic without the $O(1)$ correction of
\eqref{eq:bridge_projection_mean}. The third conditioning is where the
null becomes parameter-free.
\end{remark}

\section{The divisor-defect Gram form on orientation bits}
\label{sec:gram_form}

\begin{definition}
\label{def:hypotheses_on_G}
Let $B$ be the divisor-defect matrix on $A_N$ and $G=B^{\mathsf T}B$.
$G$ is \emph{admissible} if (G1) it is real symmetric positive
semidefinite and (G2) its kernel admits the tensor factorisation
\eqref{eq:tensor_factorisation_standing}. The orientation-bit form is
\begin{equation}
\widehat G=F^{\mathsf T}GF,\quad
\widehat G_{jk}=\sum_{a\in E_j}\sum_{b\in E_k}\Delta Q_j(a)\,G_{ab}\,\Delta Q_k(b),
\qquad
\widehat G=\sum_k\lambda_k\widehat P_k,\quad d_k=\rank\widehat P_k .
\label{eq:pushed_form}
\end{equation}
\end{definition}

\begin{proposition}
\label{prop:pushed_form_odd}
$\langle s,\widehat Gs\rangle$ is invariant under $s\mapsto-s$. Under a
single reflection $s_j\mapsto-s_j$ the terms incident to $j$,
$2\widehat G_{jk}s_js_k$ with $k\ne j$, change sign and all other terms
are unchanged. Hence the form is not $\mathcal R_N$-invariant unless
$\widehat G_{jk}=0$ for all $k\ne j$, and it is not a function of the
quadratic profile unless $\widehat G$ is diagonal.
\end{proposition}

\begin{proof}
$s_js_k$ changes sign under $s_j\mapsto-s_j$ for $k\ne j$, terms $s_ks_l$
with $k,l\ne j$ are unchanged, and $\widehat G_{jj}s_j^2$ is constant; by
Prop.~\ref{prop:reflection_fibre} functions of the profile are
$\mathcal R_N$-invariant.
\end{proof}

\section{Blindness, and its exact extension}
\label{sec:blindness}

\begin{lemma}
\label{lem:energy_as_two_point}
$\displaystyle\sum_{n\le N}M(n)^2
=\sum_{a,b\le N}\mu(a)\mu(b)\bigl(N+1-\max(a,b)\bigr).$
\end{lemma}

\begin{proof}
Expand $M(n)^2$ and exchange summation; the pair $(a,b)$ occurs for
$N+1-\max(a,b)$ values of $n$.
\end{proof}

\begin{theorem}
\label{thm:cross_excursion_vanishing}
Assume $M(N)=0$. For all $j\ne k$,
\begin{equation}
\sum_{a\in E_j}\sum_{b\in E_k}
\mu(a)\mu(b)\bigl(N+1-\max(a,b)\bigr)=0 .
\label{eq:cross_excursion_vanishing}
\end{equation}
Hence $\sum_{n\le N}M(n)^2$ depends only on within-excursion two-point
data.
\end{theorem}

\begin{proof}
For $j<k$ every $a\in E_j$ precedes every $b\in E_k$, so
$\max(a,b)=b$ and the double sum factorises as
$\bigl(\sum_{a\in E_j}\mu(a)\bigr)\bigl(\sum_{b\in E_k}\mu(b)(N+1-b)\bigr)$;
the first factor is $M(z_j)-M(z_{j-1})=0$. The case $j>k$ is symmetric.
\end{proof}

\begin{corollary}
\label{cor:profile_blindness}
No statistic factoring through $\Pi_N$ detects cross-excursion
correlation of $\mu$.
\end{corollary}

\begin{proof}
Such a statistic is $\mathcal R_N$-invariant by
Prop.~\ref{prop:reflection_fibre}; by
Thm.~\ref{thm:cross_excursion_vanishing} the cross-excursion mass does
not enter the energy, and by the action it is exactly the part on which
$\mathcal R_N$ acts nontrivially.
\end{proof}

\subsection{Two channels and the half-constant extension}
\label{subsec:two_channels}

\begin{lemma}
\label{lem:two_channels}
Let $m_j$ be an interior cut of $E_j$, $E_j^{L}=(u_j,m_j]$,
$E_j^{R}=(m_j,v_j]$, and $H_j=Q_j(m_j)$. Then
\begin{equation}
\sum_{n\in E_j^{L}}\mu_s(n)=s_jH_j,
\qquad
\sum_{n\in E_j^{R}}\mu_s(n)=-s_jH_j,
\qquad H_j\ge1 .
\label{eq:two_channels}
\end{equation}
In particular the two halves are not independent channels: they have
rank one.
\end{lemma}

\begin{proof}
Telescoping and $M(u_j)=M(v_j)=0$ give the two sums; $H_j\ge1$ because
$M$ does not vanish in the interior of an excursion.
\end{proof}

\begin{theorem}
\label{thm:half_constant_blindness}
Let $j<k$ and let a kernel $K(a,b)$ be constant in $a$ on each half of
$E_j$, say $K(a,b)=K_L(b)$ for $a\in E_j^L$ and $K_R(b)$ for
$a\in E_j^R$. Then
\begin{equation}
\sum_{a\in E_j}\sum_{b\in E_k}\mu_s(a)\mu_s(b)K(a,b)
= s_jH_j\sum_{b\in E_k}\mu_s(b)\bigl(K_L(b)-K_R(b)\bigr).
\label{eq:half_constant_blindness}
\end{equation}
If $K_L=K_R$ the exact vanishing of
Thm.~\ref{thm:cross_excursion_vanishing} is recovered; otherwise the
interaction factors through the single channel $s_jH_j$.
\end{theorem}

\begin{proof}
Split the $a$-sum over the two halves and apply
Lemma~\ref{lem:two_channels} to each.
\end{proof}

\begin{remark}
\label{rem:shift_escapes}
Theorem~\ref{thm:cross_excursion_vanishing} is the case $K_L=K_R$,
because a kernel depending on $\max(a,b)$ loses all $a$-dependence when
$j<k$. The shift kernel $K_h(a,b)=\mathbf 1[b-a=h]$ is not constant on
any half: it resolves $a=b-h$ at unit scale. Hence the shift lies
outside the extended class, and this is the precise sense in which
$C_N(h)$ escapes the excursion compression.
\end{remark}

\subsection{The maximal blind class}
\label{subsec:blind_class}

\begin{theorem}
\label{thm:blind_class}
Let $E,F$ be blocks and $P_E,P_F$ the projections onto vectors of zero
sum. For a kernel $K_{EF}$ the following are equivalent:
$x^{\mathsf T}K_{EF}y=0$ for all $x\perp\mathbf 1_E$, $y\perp\mathbf 1_F$;
$P_EK_{EF}P_F=0$; $K(a,b)=\alpha(a)+\beta(b)$ on $E\times F$. If the
block of the shift contains $r$ disjoint unit entries, its squared
Frobenius distance to this class is
\begin{equation}
r-\frac{r}{|E|}-\frac{r}{|F|}+\frac{r^2}{|E||F|}.
\label{eq:blind_distance}
\end{equation}
\end{theorem}

\begin{proof}
The first two conditions are the same statement. $P_EKP_F=0$ iff $K$
lies in the span of $\mathbf 1_E\otimes\cdot$ and $\cdot\otimes\mathbf 1_F$,
which is the additively separable class. The distance is
$\|P_EKP_F\|_F^2
=\|K\|_F^2-\frac{\|\mathbf 1^{\mathsf T}K\|^2}{|E|}
-\frac{\|K\mathbf 1\|^2}{|F|}+\frac{(\mathbf 1^{\mathsf T}K\mathbf 1)^2}{|E||F|}$,
and for $r$ disjoint unit entries these are $r$, $r/|E|$, $r/|F|$ and
$r^2/(|E||F|)$.
\end{proof}

\begin{remark}
\label{rem:blind_class_meaning}
Theorem~\ref{thm:cross_excursion_vanishing} is the case
$K(a,b)=\beta(b)$ for ordered blocks. The distance
\eqref{eq:blind_distance} is of order $r$, not a boundary error.
Globally, $U_h$ has $H$ singular values equal to $1$, so by
Eckart--Young every kernel $K_r$ of rank $r<H$ satisfies
$\|U_h-K_r\|_{\mathrm{op}}\ge1$ and
$\|U_h-K_r\|_{\mathrm{HS}}\ge\sqrt{H-r}$: no separable approximation of
sublinear rank is close in the norm needed for uniform observability.
\end{remark}

\section{The divisor--shift operator}
\label{sec:divisor_shift}

\subsection{What conjugation does not give}
\label{subsec:conjugation}

\begin{remark}
\label{rem:conjugation_nogo}
$D$ is lower triangular with $D^{-1}_{r,d}=\mu(r/d)\mathbf 1[d\mid r]$,
and
\begin{equation}
(DU_hD^{-1})_{r,m}
=\sum_{\substack{a\mid r,\ a+h\le N\\ m\mid a+h}}
\mu\!\left(\frac{a+h}{m}\right).
\label{eq:conjugated_shift}
\end{equation}
This is exact and gives no spectral gain: a similarity preserves the
problem. The entries of the commutator,
\begin{equation}
[D,U_h]_{r,m}=\mathbf 1[m>h]\mathbf 1[m-h\mid r]
-\mathbf 1[r\le N-h]\mathbf 1[m\mid r+h],
\label{eq:commutator}
\end{equation}
are not supported on primes dividing $h$: they compare the two distinct
geometries $m-h\mid r$ and $m\mid r+h$.
\end{remark}

\subsection{The difference that localises}
\label{subsec:boundary_lemma}

\begin{lemma}[Divisor--shift boundary lemma]
\label{lem:boundary}
Let $H=N-h$ and let $B_h^\circ=R_H(U_h-I)D$ be the restriction to the
interior rows $1\le n\le H$. Then
\begin{equation}
(B_h^\circ)_{n,d}=\mathbf 1[d\mid n+h]-\mathbf 1[d\mid n],
\qquad 1\le n\le H,
\label{eq:boundary_lemma}
\end{equation}
while on the full operator the rows $r>H$ equal $-\mathbf 1[d\mid r]$,
so that there the columns $d\mid h$ do not vanish. For $B_h^\circ$:
\begin{enumerate}[label=(\roman*)]
\item If $d\mid h$ the column vanishes identically.
\item If $d\nmid h$ the column is supported on the two disjoint classes
$n\equiv-h$ and $n\equiv0\pmod d$, with signs $+1$ and $-1$.
\item For every interval $I$ of integers,
\begin{equation}
\Bigl|\sum_{n\in I}\bigl(\mathbf 1[d\mid n+h]-\mathbf 1[d\mid n]\bigr)\Bigr|
\le 1 .
\label{eq:boundary_mass}
\end{equation}
\end{enumerate}
\end{lemma}

\begin{proof}
\eqref{eq:boundary_lemma} is $D_{n+h,d}-D_{n,d}$. (i) If $d\mid h$ then
$d\mid n+h\iff d\mid n$. (ii) Immediate. (iii) The sum equals the number
of multiples of $d$ in $I+h$ minus the number in $I$; two intervals of
equal length contain numbers of multiples differing by at most one.
\end{proof}

\begin{remark}
\label{rem:boundary_interpretation}
This converts the multiplicative--additive interaction into residue-class
boundaries of uniformly bounded interval discrepancy, with no regularity
assumption on $I$. The bound $1$ is optimal and improves the bound $2$
obtained by estimating the two classes separately. Version 5 stated
\eqref{eq:boundary_lemma} for the full operator; that was incorrect on
the boundary rows $r>H$ and is corrected here.
\end{remark}

\subsection{Structure and norms of the interior operator}
\label{subsec:interior_operator}

\begin{theorem}[Maximal rank]
\label{thm:interior_rank}
$\rank B_h^\circ=H=N-h$ and $\dim\ker B_h^\circ=h$, as a map
$\mathbb R^N\to\mathbb R^H$.
\end{theorem}

\begin{proof}
Take the columns $d=h+r$, $1\le r\le H$. For $n<r$ one has
$0<n+h<h+r$ and $n<h+r$, so neither $h+r\mid n+h$ nor $h+r\mid n$; for
$n=r$, $h+r\mid r+h$ and $h+r\nmid r$. The $H\times H$ submatrix on these
columns is therefore lower triangular with unit diagonal.
\end{proof}

\begin{remark}
\label{rem:no_rank_gain}
The vanishing of the columns $d\mid h$ removes only $\tau(h)$ columns and
produces no useful subspace: $B_h^\circ$ is surjective onto the interior
coordinates.
\end{remark}

\begin{proposition}[Exact norms]
\label{prop:interior_norms}
For $d\nmid h$ let $a_d=\lfloor H/d\rfloor$ and
$b_d=\lfloor N/d\rfloor-\lfloor h/d\rfloor$, the numbers of $n\le H$ with
$d\mid n$ and with $d\mid n+h$. Then the column $d$ has
$\ell^1$-norm and squared $\ell^2$-norm $a_d+b_d$, the columns $d\mid h$
vanish, and
\begin{align}
\|B_h^\circ\|_{1\to1}&=\max_{d\nmid h}(a_d+b_d),
\label{eq:norm_11}\\
\|B_h^\circ\|_{\infty\to\infty}&=\max_{n\le H}
\bigl(\tau(n)+\tau(n+h)-2\tau((n,h))\bigr),
\label{eq:norm_inf}\\
\|B_h^\circ\|_{\mathrm{HS}}^2&=\sum_{n\le H}\bigl(\tau(n)+\tau(n+h)\bigr)
-2\sum_{e\mid h}\Bigl\lfloor\frac He\Bigr\rfloor .
\label{eq:norm_hs}
\end{align}
\end{proposition}

\begin{proof}
For $d\nmid h$ the two classes are disjoint, giving the column norms.
Row $n$ is supported on divisors of exactly one of $n,n+h$; the common
divisors are those of $(n,n+h)=(n,h)$. Finally
$\sum_{n\le H}\tau((n,h))=\sum_{e\mid h}\lfloor H/e\rfloor$.
\end{proof}

\begin{proposition}[Residue-boundary form and variation bound]
\label{prop:residue_boundary}
For $d\nmid h$ write $h=q_dd+r_d$ with $0<r_d<d$ and
\begin{equation}
\mathsf A_d(t)=\Bigl\lfloor\frac{t+h}{d}\Bigr\rfloor-\Bigl\lfloor\frac hd\Bigr\rfloor
-\Bigl\lfloor\frac td\Bigr\rfloor=\mathbf 1[t\bmod d\ge d-r_d].
\label{eq:residue_block}
\end{equation}
Then $(B_h^\circ)_{t,d}=\mathsf A_d(t)-\mathsf A_d(t-1)$: every column is
the oriented boundary of a periodic block of $r_d$ consecutive residues.
Consequently, for $f$ on an interval $I=[u,v]$,
\begin{equation}
\Bigl|\sum_{n\in I}f(n)(B_h^\circ)_{n,d}\Bigr|
\le\min\bigl(|f(u)|,|f(v)|\bigr)+\operatorname{Var}_I(f),
\label{eq:variation_bound}
\end{equation}
and the endpoint term disappears when $I$ and $I+h$ contain equally many
multiples of $d$.
\end{proposition}

\begin{proof}
$\lfloor(t+h)/d\rfloor=q_d+\lfloor(t+r_d)/d\rfloor$ gives
\eqref{eq:residue_block}, and the difference in $t$ counts the new
multiple of $d$ in $(h,t+h]$ minus the new one in $(0,t]$. For
\eqref{eq:variation_bound} sum by parts from either end using the
partial sums, which have modulus at most $1$ by
Lemma~\ref{lem:boundary}(iii).
\end{proof}

\subsection{CRT pullback of the shift}
\label{subsec:crt_pullback}

\begin{lemma}
\label{lem:crt_pullback}
\begin{equation}
(D^{\mathsf T}U_hD)_{d,e}=\#\{n\le N-h: d\mid n,\ e\mid n+h\},
\label{eq:crt_pullback_definition}
\end{equation}
which vanishes unless $(d,e)\mid h$. If $(d,e)\mid h$ there is a unique
class $c\bmod[d,e]$ and
\begin{equation}
(D^{\mathsf T}U_hD)_{d,e}=\frac{N-h}{[d,e]}+\varepsilon_{d,e,h,N},
\qquad |\varepsilon_{d,e,h,N}|\le1 .
\label{eq:crt_pullback}
\end{equation}
\end{lemma}

\begin{proof}
The entry counts $n$ with $n\equiv0\pmod d$ and $n\equiv-h\pmod e$.
Compatibility of the two congruences requires
$0\equiv-h\pmod{(d,e)}$; when it holds the solution set is a single
class modulo $[d,e]$, whose count in an interval of length $N-h$ is the
stated main term with error of modulus at most $1$.
\end{proof}

\begin{remark}
\label{rem:bulk_and_boundary}
The bulk is an explicit gcd/lcm kernel and all finite irregularity sits
in a boundary term. Section~\ref{sec:haar} removes the bulk exactly.
\end{remark}

\section{Algebraic Haar on arbitrary trees, and exact bulk vanishing}
\label{sec:haar}

\begin{lemma}
\label{lem:haar_basis}
Let $I=I_0\sqcup I_1$ with masses $m_0,m_1>0$, $m=m_0+m_1$, and
\begin{equation}
\psi_I=\sqrt{\frac{m_1}{m_0m}}\,\mathbf 1_{I_0}
-\sqrt{\frac{m_0}{m_1m}}\,\mathbf 1_{I_1}.
\label{eq:haar_function}
\end{equation}
Then $\langle\psi_I,\mathbf 1_I\rangle=0$ and $\|\psi_I\|_2=1$. For any
binary tree terminating in singletons, the normalised constant together
with the $\psi_I$ form an orthonormal basis.
\end{lemma}

\begin{proof}
$\langle\psi_I,\mathbf 1_I\rangle
=\sqrt{m_0m_1/m}-\sqrt{m_0m_1/m}=0$ and
$\|\psi_I\|_2^2=m_1/m+m_0/m=1$. Orthogonality of two Haar functions
follows from disjointness of their supports or, when nested, from the
mean-zero property on the child in which the finer one is supported.
Counting gives a basis.
\end{proof}

\begin{remark}
\label{rem:no_regularity}
No doubling measure, comparable children or geometric hypothesis is
used. The identity is finite-dimensional algebra.
\end{remark}

\begin{theorem}[Exact vanishing of the CRT bulk]
\label{thm:haar_crt}
Let $R_{d,e,h}(n)=\mathbf 1[d\mid n]\mathbf 1[e\mid n+h]$. If
$(d,e)\nmid h$ then $R\equiv0$. Otherwise $R=\mathbf 1[n\equiv c\bmod q]$
with $q=[d,e]$, and writing
$\#(I_j\cap(c\bmod q))=m_j/q+\epsilon_j$ with $|\epsilon_j|<1$,
\begin{equation}
\langle R_{d,e,h},\psi_I\rangle
=\sqrt{\frac{m_1}{m_0m}}\,\epsilon_0-\sqrt{\frac{m_0}{m_1m}}\,\epsilon_1 .
\label{eq:haar_crt}
\end{equation}
Moreover the boundary defects are additive over children:
\begin{equation}
\epsilon_I(c,q)=\epsilon_{I_0}(c,q)+\epsilon_{I_1}(c,q).
\label{eq:cocycle_additivity}
\end{equation}
\end{theorem}

\begin{proof}
The compatibility statement is Lemma~\ref{lem:crt_pullback}. Insert the
counts into \eqref{eq:haar_function}: the terms $m_j/q$ contribute
$\frac1q(\sqrt{m_0m_1/m}-\sqrt{m_0m_1/m})=0$ because $\psi_I$ has mean
zero, leaving \eqref{eq:haar_crt}. Additivity holds because the counts
add and $m_I=m_{I_0}+m_{I_1}$.
\end{proof}

\begin{remark}
\label{rem:irregularity_is_data}
The irregularity of the cut is not an error term: it \emph{is} the
cocycle $(\epsilon_0,\epsilon_1)$. By \eqref{eq:cocycle_additivity} the
Haar coefficients are martingale differences of a single boundary
cocycle, which opens a telescoping across scales that the excursion
compression cannot see.
\end{remark}

\begin{proposition}[Tree-independent energy]
\label{prop:tree_independent}
For a residue class $R$ and $k_i=|R\cap I_i|$,
\begin{equation}
\langle\mathbf 1_R,\psi_I\rangle=\frac{m_1k_0-m_0k_1}{\sqrt{m_0m_1(m_0+m_1)}} .
\label{eq:haar_class}
\end{equation}
For any tree whose Haar functions complete the constant on
$\{1,\dots,H\}$ and any column $d\nmid h$ of $B_h^\circ$,
\begin{equation}
\sum_I\bigl|\langle(B_h^\circ)_{\bullet d},\psi_I\rangle\bigr|^2
=a_d+b_d-\frac{(b_d-a_d)^2}{H},
\label{eq:tree_independent}
\end{equation}
independently of the tree.
\end{proposition}

\begin{proof}
\eqref{eq:haar_class} is \eqref{eq:haar_function} evaluated on counts.
Parseval gives $\sum_I|\langle f,\psi_I\rangle|^2
=\|f\|^2-H^{-1}|\sum f|^2$; for the column,
$\|f\|^2=a_d+b_d$ and $\sum f=b_d-a_d$.
\end{proof}

\begin{warning}
\label{warn:parseval_nogo}
A tree aligned with residue boundaries can concentrate the energy in few
coefficients but cannot reduce it: irregularity produces sparsity, not
coercivity. Any new estimate must use simultaneously the
inclusion--exclusion signs, the compatibilities $(d,e)\mid h$ and the
parent--children coherence \eqref{eq:cocycle_additivity}.
\end{warning}

\begin{remark}
\label{rem:shift_matrix_in_haar}
The shift matrix in this basis is exact:
\begin{equation}
\langle\psi_I,U_h\psi_J\rangle
=\sum_{\alpha,\beta\in\{0,1\}}c_{I,\alpha}c_{J,\beta}\,
E_h(I_\alpha,J_\beta),
\qquad E_h(X,Y)=\#\{n\in X:n+h\in Y\},
\label{eq:shift_in_haar}
\end{equation}
with $c_{I,0},c_{I,1}$ the two coefficients of
\eqref{eq:haar_function}. There is an entry only when an edge joins the
supports; cancellation depends on the $2\times2$ incidence table, not on
regular lengths.
\end{remark}

\section{Observability}
\label{sec:observability}

\begin{definition}
\label{def:observability_setup}
Let $V=PH_N$, $L=D_\Lambda P|_V$ for a set $\Lambda$ of divisor rows,
and $A=PA_hP|_V$.
\end{definition}

\begin{theorem}[Exact observability criterion]
\label{thm:observability}
There exists $\kappa<\infty$ with
\begin{equation}
|x^{\mathsf T}Ax|\le\kappa\|Lx\|^2\qquad\text{for all }x\in V
\label{eq:observability_inequality}
\end{equation}
if and only if $\ker L\subseteq\ker A$. In that case
$A=L^{\mathsf T}BL$ with $B=(L^{+})^{\mathsf T}AL^{+}$ and the best
constant is
\begin{equation}
\kappa_*=\bigl\|(L^{\mathsf T}L)^{\dagger/2}A(L^{\mathsf T}L)^{\dagger/2}\bigr\|_{\mathrm{op}} .
\label{eq:best_constant}
\end{equation}
\end{theorem}

\begin{proof}
Sufficiency: write $x=x_0+x_1$ with $x_0\in\ker L$ and
$x_1\perp\ker L$. Then $Ax_0=0$ gives
$x^{\mathsf T}Ax=x_1^{\mathsf T}Ax_1$, while
$\|Lx\|=\|Lx_1\|\ge\sigma\|x_1\|$ with $\sigma$ the least positive
singular value, so \eqref{eq:observability_inequality} holds with
$\kappa=\|A\|/\sigma^2$.

Necessity: let $z\in\ker L$ and $x=z+ty$. Then
$\|Lx\|^2=t^2\|Ly\|^2$ while
$x^{\mathsf T}Ax=z^{\mathsf T}Az+2t\,z^{\mathsf T}Ay+t^2y^{\mathsf T}Ay$.
Letting $t\to0$ forces $z^{\mathsf T}Az=0$, and then the linear term must
be $O(t^2)$, forcing $z^{\mathsf T}Ay=0$ for all $y$, i.e. $Az=0$.
Formula \eqref{eq:best_constant} is the induced norm on the range.
\end{proof}

\begin{remark}
\label{rem:duality_and_sdp}
Without the kernel condition the best Frobenius fit is
$B_F=(L^{+})^{\mathsf T}AL^{+}$ with residual
$R_F=A-P_{(\ker L)^{\perp}}AP_{(\ker L)^{\perp}}$, and in operator norm
the exact dual is
\begin{equation}
\inf_B\|A-L^{\mathsf T}BL\|_{\mathrm{op}}
=\sup\bigl\{|\tr(AZ)| : Z=Z^{\mathsf T},\ \|Z\|_*\le1,\ LZL^{\mathsf T}=0\bigr\},
\label{eq:operator_norm_dual}
\end{equation}
the constraint $LZL^{\mathsf T}=0$ being exactly annihilation of the
subspace $\{L^{\mathsf T}BL\}$. For an error $\varepsilon\|x\|^2$ the
optimal pair $(\kappa,\varepsilon)$ solves the semidefinite programme
$-\kappa L^{\mathsf T}L-\varepsilon I\preceq A\preceq
\kappa L^{\mathsf T}L+\varepsilon I$.
\end{remark}

\begin{theorem}[No sublinear observability]
\label{thm:sublinear_nogo}
Fix $h\ne0$. If $q$ and $q+h$ are squarefree with $q>N/2$ and neither
row $q$ nor row $q+h$ lies in $\Lambda$, then
$x=e_q\pm e_{q+h}\in\ker L$ with $|x^{\mathsf T}A_hx|=1$ and
$\|x\|^2=2$; hence any inequality of the form
$|x^{\mathsf T}A_hx|\le\kappa\|Lx\|^2+\varepsilon\|x\|^2$ forces
$\varepsilon\ge1/2$. Consequently, to achieve $\varepsilon<1/2$ the set
$\Lambda$ must be a vertex cover of the shift graph on $(N/2,N]$, and
since that graph has maximum degree $2$ and, by Mirsky's theorem
\cite{Mirsky1949}, a positive density of edges, $|\Lambda|=\Omega_h(N)$.
No nontrivial observability is possible with $o(N)$ divisor rows.
\end{theorem}

\begin{proof}
For $q>N/2$ the only multiple of $q$ up to $N$ is $q$ itself, so column
$q$ of $D$ is seen only by row $q$; the same for $q+h$. Hence
$Lx=0$ for the stated $x$, while
$x^{\mathsf T}A_hx=\pm2\cdot\tfrac12=\pm1$ and $\|x\|^2=2$. The
covering statement follows since each uncovered edge produces such an
$x$; a graph of maximum degree $2$ with $\Omega(N)$ edges needs
$\Omega(N)$ vertices to cover.
\end{proof}

\begin{corollary}[Quantitative form]
\label{cor:sublinear_quantitative}
With $\mathfrak S_h$ the Mirsky density of \eqref{eq:mirsky_density},
any $\Lambda$ achieving $\varepsilon<1/2$ in
Theorem~\ref{thm:sublinear_nogo} satisfies
\begin{equation}
|\Lambda|\ge\frac{\mathfrak S_h}{4}\,N+o_h(N).
\label{eq:lambda_lower_bound}
\end{equation}
\end{corollary}

\begin{proof}
The shift graph on $(N/2,N]$ has $\frac{\mathfrak S_h}{2}N+o_h(N)$ edges
by \cite{Mirsky1949} and maximum degree $2$, so every vertex cover has at
least half as many vertices as there are edges.
\end{proof}

\begin{remark}
\label{rem:nogo_consequence}
Centres, quarters and returns can therefore only work after restricting
to a low-dimensional multiscale subspace; global identification is
excluded outright. The dual formula \eqref{eq:operator_norm_dual} may
equivalently be written with the constraint $QZQ=0$, $Q=L^{+}L$.
\end{remark}

\section{Centres, divisor rows and the M\"obius line}
\label{sec:centers}

\subsection{The central dichotomy}
\label{subsec:central_dichotomy}

\begin{lemma}[Central dichotomy]
\label{lem:central_dichotomy}
Let $L_j$ be any linear functional of the values or increments of $M$ on
$E_j$, fixed before the orientation is known. Then
\begin{equation}
L_j(M_s)=s_jL_j(Q_j)=s_j\kappa_j ,
\label{eq:central_dichotomy}
\end{equation}
so either $\kappa_j=0$ and $L_j$ vanishes on the whole fibre, or
$\kappa_j\ne0$ and $s_j=\operatorname{sgn}(L_j(M_s)/\kappa_j)$ is read
off directly. There is no third possibility.
\end{lemma}

\begin{proof}
Immediate from \eqref{eq:canonical_positive_path}.
\end{proof}

\begin{remark}
\label{rem:center_values_never_vanish}
With $m_j^{\mathrm{sf}}$ the support centre, i.e. the vertex after
$\ell_j$ of the $2\ell_j$ squarefree steps of $E_j$, one has
$1\le Q_j(m_j^{\mathrm{sf}})\le\ell_j$ and
$Q_j(m_j^{\mathrm{sf}})\equiv\ell_j\pmod 2$, so the central path value
never vanishes and always reads the bit. The support centre is the
canonical choice: it is squarefree, so a divisor row evaluated there
necessarily contains its own diagonal term, and it avoids spurious
plateaux with $\mu(n)=0$.
\end{remark}

\begin{warning}
\label{warn:path_reading_is_tautological}
Augmenting the Gram form with central path readings gives
$C_{\mathrm{path}}=-2\,\mathrm{diag}(h_1,\dots,h_Z)$ and hence
$G^{+}=V^{\mathsf T}AV+C_{\mathrm{path}}^{\mathsf T}WC_{\mathrm{path}}$
with $\ker G^{+}=\ker V\cap\ker C_{\mathrm{path}}=\{0\}$, since
$h_j\ne0$. The kernel disappears, but any $J_h$ then admits
$J_h=C_{\mathrm{path}}^{\mathsf T}(C_{\mathrm{path}}^{-\mathsf T}
J_hC_{\mathrm{path}}^{-1})C_{\mathrm{path}}$ with zero remainder. This is
tautological: the whole difficulty has been moved inside a diagonal bit
reader.
\end{warning}

\subsection{Sampling the divisor equation at the centre}
\label{subsec:divisor_rows_at_centers}

\begin{definition}
\label{def:cdiv}
With $m_j=m_j^{\mathrm{sf}}$ and $q_k$ the canonical positive increment
of excursion $k$, set
\begin{equation}
C^{\mathrm{div}}_{jk}=\sum_{\substack{d\mid m_j\\ d\in E_k}}q_k(d).
\label{eq:cdiv}
\end{equation}
\end{definition}

\begin{proposition}
\label{prop:cdiv_triangular}
$C^{\mathrm{div}}_{jk}=0$ for $k>j$, so $C^{\mathrm{div}}$ is lower
triangular. If moreover $u_j\ge m_j/2$, then $C^{\mathrm{div}}_{jj}
=q_j(m_j)\in\{\pm1\}$.
\end{proposition}

\begin{proof}
For $k>j$ every element of $E_k$ exceeds $m_j$, so no divisor of $m_j$
lies in $E_k$. If $u_j\ge m_j/2$ then every proper divisor of $m_j$ is
$\le m_j/2\le u_j$, hence outside $E_j=(u_j,v_j]$; only $d=m_j$ remains
and $m_j$ is squarefree.
\end{proof}

\begin{proposition}[Logarithmic exceptional set]
\label{prop:log_exceptional}
The hypothesis $u_j\ge m_j/2$ of Prop.~\ref{prop:cdiv_triangular} holds
for all but at most $\lfloor\log_2N\rfloor$ indices $j$.
\end{proposition}

\begin{proof}
Since $m_j\le v_j$, the inequality $u_j\ge v_j/2$ implies
$u_j\ge m_j/2$, so the hypothesis can fail only when $v_j>2u_j$. Since
$M(1)=1$ and $M(2)=0$, $z_1=2$; thus $u_1=0$ and $u_j\ge2$ for $j\ge2$.
Let $2\le j_1<\dots<j_r$ be the failing indices with $j\ge2$. Using
$u_{j+1}=v_j$ and monotonicity of $u$,
$N\ge v_{j_r}>2u_{j_r}>2^2u_{j_{r-1}}>\dots>2^ru_{j_1}\ge2^{r+1}$,
so $r+1\le\lfloor\log_2N\rfloor$. Adding the possible failure at $j=1$
gives at most $\lfloor\log_2N\rfloor$ failing indices.
\end{proof}

\begin{theorem}
\label{thm:mobius_line}
Let $\mathcal J$ be the set of indices satisfying $u_j\ge m_j/2$. In
orientation coordinates, treating the row $n=1$ affinely by $w(1)=1$,
\begin{equation}
C^{\mathrm{div}}s_\mu=0,
\qquad
\rank C^{\mathrm{div}}\ \ge\ |\mathcal J|\ \ge\ Z-\lfloor\log_2N\rfloor .
\label{eq:mobius_line}
\end{equation}
In particular $\dim\ker C^{\mathrm{div}}\le\lfloor\log_2N\rfloor+1$ and
$s_\mu$ lies in that kernel.
\end{theorem}

\begin{proof}
Row $j$ of $C^{\mathrm{div}}s_\mu$ equals $\sum_{d\mid m_j}\mu(d)$,
which vanishes for $m_j>1$; the affine treatment of $d=1$ fixes the
remaining normalisation. The rank bound follows from lower
triangularity together with Prop.~\ref{prop:cdiv_triangular} on
$\mathcal J$ and Prop.~\ref{prop:log_exceptional}.
\end{proof}

\begin{corollary}[Compressed selection]
\label{cor:selection_count}
The number of $s\in\{\pm1\}^Z$ satisfying the central divisor equations,
with the affine normalisation of the row $n=1$, is at most
$2^{\lfloor\log_2N\rfloor+1}\le2N$.
\end{corollary}

\begin{proof}
By Theorem~\ref{thm:mobius_line} the solution set is an affine subspace
of dimension $k\le\lfloor\log_2N\rfloor+1$. It projects injectively onto
some $k$ coordinates, so its points with entries $\pm1$ are determined by
those $k$ entries, giving at most $2^k$ of them.
\end{proof}

\begin{corollary}[Completion to unique selection]
\label{cor:selection_completion}
Assume $M(N)=0$. With the canonical $F$ of \eqref{eq:canonical_F},
$C^{\mathrm{div}}_{jk}=(DF)_{m_j,k}$, so $C^{\mathrm{div}}$ consists of rows
of $DF$, which has rank $Z$. Hence at most $\lfloor\log_2N\rfloor+1$
further rows of $DF$ complete $C^{\mathrm{div}}$ to a system of rank $Z$,
whose unique solution with the normalisation of the row $n=1$ is $s_\mu$.
\end{corollary}

\begin{proof}
$(DF)_{r,k}=\sum_{d\mid r}F_{d,k}$, and for $r=m_j$ this is
\eqref{eq:cdiv}. $D$ is unitriangular and $F$ has $Z$ nonzero columns with
disjoint supports, so $\rank DF=Z$; extend a basis of the row space of
$C^{\mathrm{div}}$ to one of $DF$, which needs at most
$Z-\rank C^{\mathrm{div}}\le\lfloor\log_2N\rfloor+1$ rows. Since
$DFs_\mu=D\mu=e_1$, every row $r\ne1$ vanishes at $s_\mu$ and the row
$r=1$ fixes the normalisation.
\end{proof}

\begin{remark}
\label{rem:selection_meaning}
This is a near-minimal compression of M\"obius inversion inside the
fibre: $Z+O(\log N)$ divisor equations instead of $N$ select the
M\"obius orientation among $2^Z$. It works inside a fibre that already
fixes $|M(n)|$ for every $n$, so it carries no information about the
size of $M$.
\end{remark}

\begin{remark}
\label{rem:cdiv_meaning}
This is a sparse orientational selector: roughly one divisor equation per
excursion, chosen canonically at its centre, isolating the M\"obius line
up to a kernel of logarithmic dimension. For the full divisor operator
these rows are already present, $C^{\mathrm{div}}=RV$ with $R$ a row
selector, so they do not reduce $\ker V$; for a truncation $V_y$ they
reduce the kernel exactly when
$C^{\mathrm{div}}|_{\ker V_y}\ne0$.
\end{remark}

\begin{warning}
\label{warn:isolating_is_not_bounding}
Isolating the M\"obius line does not bound the shift on it. If the
central rows reduce the kernel to $\operatorname{span}\{s_\mu\}$, then
estimating $J_h$ on that kernel is exactly estimating the
cross-excursion M\"obius correlation. Centres help as an arithmetic
selector, not yet as an estimator.
\end{warning}

\subsection{Conditioning the two-point expansion on central data}
\label{subsec:conditioning}

\begin{proposition}
\label{prop:walsh_conditioning}
Write the expansion \eqref{eq:fibre_expansion} in the orientation bits
as
\begin{equation}
C_N(h;s)=A_{N,h}+\ell_h^{\mathsf T}s+s^{\mathsf T}J_h^{\mathrm{off}}s ,
\label{eq:walsh_expansion}
\end{equation}
and split $s=(a,u)$ with $a\in\{\pm1\}^{F}$ fixed and $u$ free. Then
\begin{equation}
C_N(h;s)=A'+\bigl(\ell_U+2J_{UF}\,a\bigr)^{\mathsf T}u+u^{\mathsf T}J_{UU}\,u,
\qquad
A'=A_{N,h}+\ell_F^{\mathsf T}a+a^{\mathsf T}J_{FF}\,a .
\label{eq:walsh_conditioning}
\end{equation}
\end{proposition}

\begin{proof}
Write $J_h^{\mathrm{off}}$ in blocks indexed by $F$ and $U$ and use its
symmetry, so that the two mixed blocks contribute $2a^{\mathsf T}J_{FU}u$.
\end{proof}

\begin{remark}
\label{rem:conditioning_gives_nothing}
Conditioning on all signed central values fixes $s$, by
Lemma~\ref{lem:central_dichotomy}, and the fibre is a singleton: every
Walsh coefficient disappears only because the full vector has been
supplied. Conditioning on central magnitudes $|M(m_j)|$ changes nothing,
since they are profile data. Conditioning on a proper subset $F$
converts fixed--free quadratic terms into linear ones but leaves the
quadratic form $J_{UU}$ on the unobserved bits intact. No quantitative
reduction follows without an independent bound on $J_{UU}$.
\end{remark}

\begin{remark}
\label{rem:quarters_octaves}
On orientations, a single nonvanishing odd evaluation already reads
$s_j$, so quarter and octave evaluations are redundant by
Lemma~\ref{lem:central_dichotomy}. On the shift, a coarse Haar
truncation does not approximate $A_h$ uniformly on vectors with entries
in $\{-1,0,1\}$, which may oscillate at unit scale; deterministic
control of the remainder requires cells of size $O(h)$, that is,
essentially full resolution for fixed $h$.
Proposition~\ref{prop:leaf_energy} makes this exact.
\end{remark}

\section{Intrablock identities}
\label{sec:intrablock}

\begin{lemma}
\label{lem:second_difference}
Let $x_1,\dots,x_L$ have zero sum, $P_t=\sum_{r\le t}x_r$ extended by
zero outside $[0,L]$, and $A_h=\sum_tP_tP_{t+h}$. Then
\begin{equation}
C_h:=\sum_tx_tx_{t+h}=2A_h-A_{h-1}-A_{h+1}.
\label{eq:second_difference}
\end{equation}
Moreover
\begin{equation}
C_0+2\sum_{h\ge1}C_h=0,
\qquad
\sum_{t=1}^{L-1}P_t^2=-\sum_{h\ge1}hC_h .
\label{eq:aggregate_identities}
\end{equation}
\end{lemma}

\begin{proof}
Expand $x_t=P_t-P_{t-1}$ and collect the four autocorrelations. The
first identity in \eqref{eq:aggregate_identities} is
$\sum_{h\in\mathbb Z}C_h=(\sum_tx_t)^2=0$. For the second, apply
\eqref{eq:second_difference} for all $h\in\mathbb Z$ and reindex:
$\sum_h|h|C_h=\sum_hA_h(2|h|-|h+1|-|h-1|)=-2A_0$, since the bracket
vanishes for $h\ne0$ and equals $-2$ at $h=0$; halve.
\end{proof}

\begin{remark}
\label{rem:intrablock_is_profile}
Inside an excursion $P_tP_{t+h}=|P_t||P_{t+h}|$, so the intrablock
autocorrelation is already determined by the absolute profile. The
centre adds nothing there.
\end{remark}

\begin{remark}[The centre does not determine a lag]
\label{rem:counterexample_lag}
The two positive excursions of length $8$
\begin{equation}
(+,+,+,-,+,-,-,-)\quad\text{and}\quad(+,+,-,+,-,+,-,-)
\label{eq:counterexample_words}
\end{equation}
have paths $(1,2,3,2,3,2,1,0)$ and $(1,2,1,2,1,2,1,0)$, hence the same
central value $2$, but $C_1=1$ and $C_1=-3$ respectively. Length, zero
sum and central height do not control even the first lag.
\end{remark}

\begin{lemma}[Left--right splitting]
\label{lem:left_right}
With $m$ the centre and
$C_h=C_h^{LL}+C_h^{RR}+C_h^{\times}$, only indices
$m-h<t\le m$ can cross, so $|C_h^{\times}|\le h$ and, summing over
excursions, $|\sum_jC_{h,j}^{\times}|\le hZ$.
\end{lemma}

\begin{proof}
A pair $(t,t+h)$ crosses $m$ iff $t\le m<t+h$.
\end{proof}

\begin{remark}
\label{rem:Z_growth_is_the_hypothesis}
Lemma~\ref{lem:left_right} is $o(N)$ precisely when $Z=o(N)$, which is
not proved anywhere in this project; see \ref{item:Q5}. The other two
parts can be of the size of the excursion length: the centre isolates a
layer of width $h$ only.
\end{remark}

\begin{lemma}[Discrete coarea]
\label{lem:coarea}
\begin{equation}
\mu(n)=\sum_{\ell\in\mathbb Z}
\Bigl(\mathbf 1[M(n-1)\le\ell<M(n)]-\mathbf 1[M(n)\le\ell<M(n-1)]\Bigr).
\label{eq:coarea}
\end{equation}
Consequently $C_N(h)$ is a sum of oriented intersections of level-set
boundaries separated by $h$.
\end{lemma}

\begin{proof}
If $\mu(n)=+1$ the first indicator contributes once and the second is
empty; if $\mu(n)=-1$ the roles are exchanged; if $\mu(n)=0$ both are
empty.
\end{proof}

\begin{proposition}[Global second difference]
\label{prop:global_second_difference}
With $H=N-h$ and $A_k(L)=\sum_{n=1}^{L}M(n)M(n+k)$,
\begin{equation}
C_N(h)=2A_h(H-1)-A_{h-1}(H-1)-A_{h+1}(H-1)+M(H)\mu(N).
\label{eq:global_second_difference}
\end{equation}
For a partition into $B$ blocks of zero sum, i.e. between returns of $M$
to a common level, not necessarily zero, with local prefixes $P_j$
extended by zero and $A_j(k)=\sum_nP_j(n)P_j(n+k)$,
\begin{equation}
C_N(h)=\sum_j\bigl(2A_j(h)-A_j(h-1)-A_j(h+1)\bigr)+R_\partial+R_{\mathrm{tail}},
\qquad |R_\partial|\le h(B-1),
\label{eq:block_second_difference}
\end{equation}
with $|R_{\mathrm{tail}}|$ bounded by the length of the terminal tail.
\end{proposition}

\begin{proof}
Expand $\mu(n)=M(n)-M(n-1)$ in $\sum_{n\le H}\mu(n)\mu(n+h)$: the four
sums are $A_h(H)$, $A_{h-1}(H)$, $A_{h+1}(H-1)$ and $A_h(H-1)$, and
$A_h(H)-A_{h-1}(H)=A_h(H-1)-A_{h-1}(H-1)+M(H)\mu(N)$. The block version is
Lemma~\ref{lem:second_difference} on each block, with at most $h$ crossing
pairs per cut.
\end{proof}

\begin{remark}
\label{rem:curvature_is_the_content}
The uncontrolled term is the curvature at scale $h$ of the local path
autocorrelations. Centres and hitting times can shorten the tail, but
each new cut costs up to $h$ pairs, and no universal choice makes both
the number of cuts and the curvatures small. In the coarea form
$C_N(h)=\sum_{\ell,m}\sum_{n\le H}\sigma_\ell(n)\sigma_m(n+h)$,
up- and down-crossings balance globally in closed blocks, but their
shifted products do not pair: the shift selects particular times, not
totals.
\end{remark}

\section{Exact local structure at each prime}
\label{sec:local}

\begin{proposition}
\label{prop:local_tables}
Fix a prime $p$ and $t=v_p(h)$. Counting residues modulo $p^2$, the
admissible states $(v_p(n),v_p(n+h))\in\{0,1\}^2$ and their
multiplicities are
\begin{center}
\begin{tabular}{c|l|c}
$t$ & states and multiplicities & local sign\\ \hline
$t=0$ & $(0,0)$: $p^2-2p$; $(1,0)$, $(0,1)$: $p-1$ each & $+,-,-$\\
$t=1$ & $(0,0)$: $p^2-p$; $(1,1)$: $p-2$ & $+$\\
$t\ge2$ & $(0,0)$: $p^2-p$; $(1,1)$: $p-1$ & $+$
\end{tabular}
\end{center}
The local squarefree density of the pair is $1-2/p^2$ if $p^2\nmid h$
and $1-1/p^2$ if $p^2\mid h$, which is the Mirsky structure
\cite{Mirsky1949}. For any finite set of primes $P$,
\begin{equation}
\prod_{p\in P}(-1)^{v_p(n)+v_p(n+h)}
\mathbf 1[v_p(n),v_p(n+h)\le1]
\label{eq:local_identity}
\end{equation}
is an exact identity, not an independence hypothesis.
\end{proposition}

\begin{proof}
For $p\nmid h$ the classes $n\equiv0$ and $n\equiv-h\pmod p$ are
distinct, giving $(p-2)p$ residues with $(0,0)$ and $p-1$ each with
$(1,0)$, $(0,1)$; the two excluded residues are $p^2\mid n$ and
$p^2\mid n+h$. For $p\|h$, $p\mid n$ forces $p\mid n+h$, and among the
$p$ residues with $p\mid n$ one must delete $n\equiv0$ and
$n\equiv-h$, distinct modulo $p^2$, leaving $p-2$. For $p^2\mid h$ the
two deleted residues coincide, leaving $p-1$. Signs are
$(-1)^{v_p(n)+v_p(n+h)}$; the densities follow by counting the excluded
residues. CRT makes the joint count over distinct primes multiplicative.
\end{proof}

\begin{remark}
\label{rem:local_is_not_the_obstruction}
At the level of residues the shift does not break the local structure:
every prime participates, with a universal off-diagonal factor for
$p\nmid h$ and a diagonal factor for $p\mid h$, and CRT is exact. What
fails on globalising is the cut $n\le N$, the infinite product and the
parity of large primes appearing to the first power.
Proposition~\ref{prop:valuation_nogo} shows that at the level of
valuations the shift does break it.
\end{remark}

\begin{proposition}[Signed local sums and Mirsky density]
\label{prop:signed_local}
With $w_p(n)=(-1)^{\mathbf 1[p\mid n]+\mathbf 1[p\mid n+h]}
\mathbf 1[p^2\nmid n]\mathbf 1[p^2\nmid n+h]$, the sum of $w_p(n)$ over
$n\bmod p^2$ equals $p^2-4p+2$, $p^2-2$ or $p^2-1$ according as
$v_p(h)=0$, $1$ or $\ge2$. With $q(n)=\sum_{r^2\mid n}\mu(r)$,
\begin{equation}
\sum_{n\le H}q(n)q(n+h)
=H\!\!\sum_{\substack{r,s\\ (r,s)^2\mid h}}\!\!\frac{\mu(r)\mu(s)}{[r^2,s^2]}
+\sum_{\substack{r,s\\ (r,s)^2\mid h}}\mu(r)\mu(s)\,\epsilon_{r,s},
\qquad|\epsilon_{r,s}|<1,
\label{eq:mirsky_expansion}
\end{equation}
the sums over $r\le\sqrt H$, $s\le\sqrt{H+h}$, and the main term has
density
\begin{equation}
\mathfrak S_h=\prod_{p^2\mid h}\Bigl(1-\frac1{p^2}\Bigr)
\prod_{p^2\nmid h}\Bigl(1-\frac2{p^2}\Bigr)>0 .
\label{eq:mirsky_density}
\end{equation}
\end{proposition}

\begin{proof}
The signed sums follow from the table of
Proposition~\ref{prop:local_tables}. For \eqref{eq:mirsky_expansion},
expand both $q$ factors; the system $r^2\mid n$, $s^2\mid n+h$ is
solvable iff $(r,s)^2\mid h$, and then counts one class modulo
$[r^2,s^2]$. The density is the product of the local densities of
Proposition~\ref{prop:local_tables}.
\end{proof}

\begin{remark}
\label{rem:mirsky_error}
Mirsky \cite{Mirsky1949} controls the error in \eqref{eq:mirsky_expansion}
by a power saving, recalled as $O_h(H^{2/3+o(1)})$ and not verified here.
The tail $\sum_{n\le N}|q(n)-q_P(n)|\le\sum_{P<p\le\sqrt N}\lfloor N/p^2\rfloor$
is controlled, but it controls only large squares, not the parity of
large primes appearing to the first power.
\end{remark}

\begin{proposition}[The shift does not descend to valuations]
\label{prop:valuation_nogo}
Let $p\nmid h$ and choose units $u,v\bmod p$ with $u\equiv-h$ and
$v\not\equiv-h\pmod p$. Then $v_p(u)=v_p(v)=0$, while
$v_p(u+h)\ge1$ and $v_p(v+h)=0$. Hence translation by $h$ is not well
defined on the quotient remembering only $v_p(n)$, and there is no exact
intertwiner between valuation-indexed local factors such as the $L_p$ of
\ref{item:S4} and the shift without adjoining residue information
$n/p^{v_p(n)}\bmod p^k$.
\end{proposition}

\begin{proof}
Immediate.
\end{proof}

\section{Multiplicative weights collapse}
\label{sec:collapse}

\begin{theorem}
\label{thm:multiplicative_collapse}
Let $m_d(x)=\sum_{n\le x,\ (n,d)=1}\mu(n)/n$. For $p\nmid d$,
\begin{equation}
m_d(x)=m_{dp}(x)-\frac1p\,m_{dp}(x/p),
\qquad
m_{dp}(x)=\sum_{k\ge0}p^{-k}m_d(x/p^k),
\label{eq:local_action}
\end{equation}
the second sum being finite. Consequently
\begin{equation}
m_d(x)=\sum_{\substack{e\ge1\\ \rad(e)\mid d}}\frac1e\,m_1(x/e),
\label{eq:lattice_closes}
\end{equation}
and with $G_{ab}=2N/[a,b]$,
\begin{equation}
\frac{\langle\mu,G\mu\rangle}{2N}
=\sum_{d\ \mathrm{squarefree}}\varphi(d)
\left(\sum_{\rad(u)=d}\frac{m_1(N/u)}{u}\right)^{2}.
\label{eq:gram_collapse}
\end{equation}
In Mellin variables the local operator has symbol $1-p^{-1-s}$ and the
full product has symbol $1/\zeta(1+s)$.
\end{theorem}

\begin{proof}
Split $\{(n,d)=1\}$ according to $p\mid n$; for $n=pn'$ squarefree with
$(n',dp)=1$ one has $\mu(n)=-\mu(n')$, giving the first identity. The
Neumann series inverts it and terminates because $m_d(x/p^k)=0$ for
$p^k>x$. Composing over $p\mid d$ gives \eqref{eq:lattice_closes}.
For \eqref{eq:gram_collapse}, write $\gcd(a,b)=\sum_{d\mid a,\ d\mid b}\varphi(d)$
to obtain
$\langle\mu,G\mu\rangle/2N=\sum_d\varphi(d)d^{-2}\mu(d)^2m_d(N/d)^2$;
substituting \eqref{eq:lattice_closes} with $u=de$, the conditions
$\rad(e)\mid d$, $d\mid u$ and $d$ squarefree force $d=\rad(u)$, so each
$u$ lies in exactly one block and $1/e=d/u$. The Mellin symbol follows
from $\sigma_pf(x)=f(x/p)\leftrightarrow p^{-s}$.
\end{proof}

\begin{corollary}
\label{cor:gcd_weight_is_inert}
Expanding $G$ along its diagonals and using
$\gcd(a,a+h)=\gcd(a,h)$,
\begin{equation}
2\sum_{h\ge1}\sum_a\mu(a)\mu(a+h)\frac{\gcd(a,h)}{a(a+h)}
=\sum_{d}\varphi(d)\Bigl(\sum_{\rad(u)=d}\frac{m_1(N/u)}{u}\Bigr)^2
-\sum_{a\le N}\frac{\mu(a)^2}{a^2},
\label{eq:gcd_aggregate}
\end{equation}
so the gcd-weighted Chowla aggregate is determined by $m_1$, i.e. by
$1/\zeta(1+s)$. The left side is in particular bounded below by
$-\zeta(2)/\zeta(4)$ since the right side is a sum of squares.
\end{corollary}

\begin{proof}
Immediate from Theorem~\ref{thm:multiplicative_collapse} and
$G_{a,a+h}=2N\gcd(a,h)/(a(a+h))$.
\end{proof}

\begin{warning}[Scope of the collapse]
\label{warn:dichotomy}
Theorem~\ref{thm:multiplicative_collapse} proves collapse for the
gcd/lcm kernel. The mechanism, closure of the divisor lattice, suggests
that diagonal weights with a multiplicative symbol collapse in the same
way; this is heuristic and is not proved here in general. The claim made
in Version 4, that every weight without a multiplicative symbol is
strictly harder than $C_N(h)$ with no intermediate case, is incorrect in
that generality and is withdrawn. The proof of
collapse uses $\mu(da')=\mu(d)\mu(a')$ and uniqueness of the radical,
i.e. commutativity of the Euler product; a non-commutative deformation
such as the cocycle $\Phi(\pi)$ of \ref{item:S7} is the only place where
the argument does not apply. That escape is a fixed jet in $\vartheta$,
and \ref{item:S8} closes fixed jets; only the scaling window
$\vartheta\asymp\vartheta_c(N)$ survives.
\end{warning}

\section{Scale-resolved statistic and sine-kernel content}
\label{sec:scale_resolved}

\begin{definition}
\label{def:scale_resolved_statistic}
$\mathcal E_k=\langle s,\widehat P_ks\rangle$,
$\mathcal R_k=\mathcal E_k/d_k$.
\end{definition}

\begin{proposition}
\label{prop:statistic_calibration}
Under $\nu$, $\mathbb E_\nu[\mathcal R_k]=1$ exactly and
$\Var_\nu[\mathcal R_k]
=2d_k^{-2}\bigl(d_k-\sum_j\widehat P_{k,jj}^2\bigr)$, with
$\sum_k\mathcal E_k=Z=\sum_kd_k$.
\end{proposition}

\begin{proof}
Prop.~\ref{prop:exact_isotropy} with $A=\widehat P_k$, and
$\sum_k\widehat P_k=I_Z$.
\end{proof}

\begin{lemma}
\label{lem:band_projector_sinc}
The orthogonal projection onto a contiguous block of $d$ characters of a
cyclic group of order $Z$ has Dirichlet kernel, whose scaling limit with
$d/Z$ fixed is $\sin(\pi x)/(\pi x)$ \cite{Mehta2004}. Equivalently
\begin{equation}
1-R_2(x)=\Bigl(\frac{\sin\pi x}{\pi x}\Bigr)^{2}
=\int_{-1}^{1}(1-|\tau|)e^{2\pi i\tau x}\,d\tau .
\label{eq:fejer_identity}
\end{equation}
\end{lemma}

\begin{warning}
\label{warn:sinc_is_the_null}
By Prop.~\ref{prop:exact_isotropy} the reflection null has covariance
$I$, which restricted to a band of characters is a sine kernel. The sine
kernel is therefore the \emph{null} of $\mathcal R_k$, not a consequence
of it, and deriving it from flatness carries no evidential weight about
the zeros of $\zeta$.
\end{warning}

\begin{remark}
\label{rem:multiplicative_band}
Hypothesis (G2), i.e. \ref{item:S4}, makes the eigenbasis of $G$ a
character basis of the multiplicative structure, so the transported band
gives a sine kernel in $\log n$, the variable in which
$n^{i\gamma}=e^{i\gamma\log n}$ oscillates. This does not affect
Warning~\ref{warn:sinc_is_the_null}.
\end{remark}

\begin{remark}
\label{rem:nontautological_criterion}
A statistic carries sine-kernel content non-tautologically if and only if
its null value is a number different from its sine-kernel value; the
measured position between the two is the coefficient. The three
quantities of Version 2 satisfying this are retained below.
\end{remark}

\begin{definition}
\label{def:superposition}
For the superposition of $p$ independent identically distributed spectra
each of density $1/p$,
\begin{equation}
R_2^{(p)}(x)=1-\frac1p\Bigl(\frac{\sin(\pi x/p)}{\pi x/p}\Bigr)^{2},
\qquad
Y^{(p)}_{jk}=\frac1p\Bigl(\frac{\sin(\pi(j-k)/p)}{\pi(j-k)/p}\Bigr)^{2}.
\label{eq:superposition_matrix}
\end{equation}
\end{definition}

\begin{proposition}
\label{prop:positive_sine_fraction}
If $\tilde r$ is strictly monotone in $p$ with $\tilde r(\infty)=2\ln2-1$
and $\tilde r(1)$ the unitary value, then a measured
$\tilde r>2\ln2-1$ forces $p<\infty$ and hence a strictly positive
sine-kernel fraction $1/p$.
\end{proposition}

\begin{proof}
Strict monotonicity and the endpoint values.
\end{proof}

\begin{lemma}
\label{lem:triangle_moments}
$\displaystyle m_n=\int_0^1(1-\tau)\tau^n\,d\tau=\frac{1}{(n+1)(n+2)}$,
and $(1-\tau)_+$, having compact support, is determined by
$\{m_n\}_{n\ge0}$.
\end{lemma}

\begin{proof}
Elementary integration and determinacy of the Hausdorff moment problem.
\end{proof}

\begin{lemma}
\label{lem:spacing_variance_anchors}
Under the Wigner surmises,
$\Var_{\mathrm{Poisson}}(s)=1$,
$\Var_{\mathrm{GOE}}(s)=4/\pi-1$,
$\Var_{\mathrm{GUE}}(s)=3\pi/8-1$.
\end{lemma}

\begin{proof}
For $p(s)=\tfrac\pi2se^{-\pi s^2/4}$ and
$p(s)=\tfrac{32}{\pi^2}s^2e^{-4s^2/\pi}$, both normalised to
$\langle s\rangle=1$, ratios of Gaussian moments give
$\langle s^2\rangle=4/\pi$ and $3\pi/8$; the Poisson case is the
exponential law.
\end{proof}

\begin{remark}
\label{rem:empirical_inputs}
Recorded in separate work, not computed here, and used only to indicate
the scale of the effects: a spacing ratio $\tilde r\approx0.42$ for the
spectrum of the smooth arithmetic term; relative freeness defects of
about $12\%$ at order two and $32\%$ at order three; an energy deficit
ratio of about $0.137$ for the area energy and about $0.052$ for the
truncated logarithmic energy; a persistent bias of about $2.5\%$; and the
systematic observation that $d=1$ is the weakest anomaly. The reference
values $2\ln2-1$, $\approx0.5307$ and $\approx0.5996$ are recalled
from \cite{AtasEtAl2013} and not verified here.
\end{remark}

\begin{remark}
\label{rem:defect_ratio_discrimination}
The triangle moments give $m_3/m_2=0.6$, decreasing; Chebyshev-basis
variance of linear statistics grows linearly \cite{Johansson1998},
giving $1.5$; the recorded defects give $32/12\approx2.67$. The observed
growth is faster than linear where the triangle requires decay. The
ratio is invariant under a global rescaling of the normalisation of
relative defect, so the discrimination survives that convention.
\end{remark}

\begin{remark}
\label{rem:symmetry_class_obstruction}
$\mathcal R_N$ and $\widehat G$ are real, hence $\beta=1$, whose form
factor $2\tau-\tau\log(1+2\tau)$ on $\tau\le1$ \cite{Mehta2004} is not a
band indicator; the conjectural class for the zeros is $\beta=2$
\cite{Montgomery1973, Odlyzko1987}.
\end{remark}

\begin{remark}[Withdrawn conjecture]
\label{conj:twisted_orientation_bits}
Versions 1--4 conjectured that letting $U(1)^Z$ act by
$s_j\mapsto e^{i\vartheta_j}s_j$ and transporting $\widehat G$ produces an
orthogonal-to-unitary crossover at $\vartheta_c\asymp1/Z$. As stated this
is incorrect: by Proposition~\ref{prop:gauge} a transported form
$U_\vartheta^{*}\widehat GU_\vartheta$ is isospectral for every
$\vartheta$, and the scale $1/Z$ does not follow from
\cite{PandeyMehta1983}. See Question~\ref{q:flux_deformation}.
\end{remark}

\section{Audit results of Version 5}
\label{sec:audit}

\subsection{The energy kernel sits at $\alpha=0$}
\label{subsec:alpha_zero}

\begin{proposition}
\label{prop:alpha_localisation}
Let $K(\Delta)=1/(2+i\Delta)$ with $\Delta=\gamma-\gamma'$. Then
\begin{equation}
K(\Delta)+K(-\Delta)=\frac{4}{4+\Delta^2}
=\int_{-\infty}^{\infty}e^{-2|u|}e^{-iu\Delta}\,du .
\label{eq:montgomery_weight}
\end{equation}
Under $u=\alpha\log T$ the induced normalised density on $\alpha$ is
\begin{equation}
p_T(\alpha)=\log T\;T^{-2|\alpha|},
\qquad
\mathbb E|\alpha|=\frac{1}{2\log T},
\qquad
\Pr(|\alpha|>A)=T^{-2A}.
\label{eq:alpha_density}
\end{equation}
\end{proposition}

\begin{proof}
Sum the two fractions; $\int e^{-a|u|}e^{-iu\Delta}du=2a/(a^2+\Delta^2)$
with $a=2$. Change variables; $\int p_T=1$, and the moments are
elementary exponential integrals.
\end{proof}

\begin{remark}
\label{rem:alpha_meaning}
The hermitised kernel is Montgomery's smoothing weight
\cite{Montgomery1973} without the factor $T^{i\alpha(\gamma-\gamma')}$,
so the scalar energy tests the pair correlation at $\alpha=0$, smeared
over width $1/\log T$. Evaluating the closed forms at the height of the
$400$th zero, $T\simeq679.74$, gives $\mathbb E|\alpha|\approx0.0767$ and
$\Pr(|\alpha|>1)\approx2.16\times10^{-6}$. To reach $\alpha_0>1$ one needs
the modulated family
\begin{equation}
K_{\alpha_0,T}(\Delta)=\frac{4}{4+\Delta^2}\,e^{i\alpha_0\log T\,\Delta}.
\label{eq:modulated_kernel}
\end{equation}
The dependence of the energy on $\log N$
(Remark~\ref{rem:negative_result_energy_scalar}) supplies exactly this
phase at $\log N=\alpha_0\log T$; what it does not supply is localisation
in height (Remark~\ref{rem:negative_result_height_weighting}). The
modulated family is available, the localised one is not.
\end{remark}

\begin{warning}
\label{warn:two_alphas}
The Montgomery variable $\alpha$ of this subsection is unrelated to the
parameter $\alpha$ of the ordinal-prime programme behind \ref{item:S9},
for which separate work records $\sigma_c(\alpha)=1+\Theta$ and
$\mathrm{RH}\iff\sigma_c(\alpha)=3/2$. The two must not be conflated.
\end{warning}

\subsection{Gauge obstruction to the crossover}
\label{subsec:gauge}

\begin{proposition}
\label{prop:gauge}
\begin{enumerate}[label=(\roman*)]
\item A Hermitian matrix whose off-diagonal support is the edge set of a
forest is conjugate, by a diagonal unitary, to a real symmetric matrix
with nonnegative off-diagonal entries. In particular a Hermitian Jacobi
matrix with entries $a_je^{\pm i\phi_j}$ has the spectrum of the real
Jacobi matrix with entries $a_j$.
\item If $\widehat G^{(\vartheta)}=U_\vartheta^{*}\widehat GU_\vartheta$
with $U_\vartheta$ unitary, then
$\operatorname{Spec}\widehat G^{(\vartheta)}=\operatorname{Spec}\widehat G$
for every $\vartheta$.
\end{enumerate}
\end{proposition}

\begin{proof}
(i) Root each tree and choose the phase of each vertex recursively along
edges so that every edge entry becomes real and nonnegative; no cycle
imposes a compatibility condition. (ii) Unitary similarity.
\end{proof}

\begin{question}
\label{q:flux_deformation}
Construct a Hermitian deformation of $\widehat G$, or of a divisibility
operator, on a graph with cycles, with phases whose fluxes around cycles
are not coboundaries, for instance built from the cocycle $\Phi(\pi)$ of
\ref{item:S7}. Only such a deformation can change the spectrum.
\end{question}

\begin{heuristic}
\label{heur:crossover_scale}
For a non-isospectral deformation $H_\vartheta$ with
$A=\frac{d}{d\vartheta}H_\vartheta\big|_{\vartheta=0}$, the crossover
scale is expected at
\begin{equation}
\vartheta_c\asymp
\frac{\Delta_Z}{\sqrt{\mathbb E|\langle u_j,Au_k\rangle|^2}},
\label{eq:crossover_scale}
\end{equation}
$\Delta_Z$ the mean level spacing and $u_j$ the eigenvectors, in the
spirit of \cite{PandeyMehta1983}. Depending on normalisation this may be
$Z^{-1}$, $Z^{-1/2}$ or another scale; the coefficients $A_r,B_r$ of
\ref{item:S7} alone do not decide it.
\end{heuristic}

\subsection{Hermitian collision versus the M\"obius energy constant}
\label{subsec:cherm_energy}

\begin{proposition}
\label{prop:cherm_energy}
Assume RH and simple zeros, and let
$\beta_M=\sum_\rho|\rho\zeta'(\rho)|^{-2}
=2\sum_{\gamma>0}|\rho|^{-2}|\zeta'(\rho)|^{-2}$. For every
$\gamma>0$,
\begin{equation}
\frac{2}{|\rho|^2|\zeta'(\rho)|^2}-C_{\mathrm{herm}}(\rho)
=\frac{1}{|\rho|^4|\zeta'(\rho)|^2},
\label{eq:cherm_pointwise}
\end{equation}
hence, whenever the sums converge,
\begin{equation}
\sum_{\gamma>0}C_{\mathrm{herm}}(\rho)
=\beta_M-\sum_{\gamma>0}\frac{1}{|\rho|^4|\zeta'(\rho)|^2}.
\label{eq:cherm_sum}
\end{equation}
\end{proposition}

\begin{proof}
With $\beta=\tfrac12$, $|\rho|^2=\gamma^2+\tfrac14$ and
$\gamma^2-\beta^2=\gamma^2-\tfrac14$, so
$\frac{2}{|\zeta'|^2}\bigl(\frac{1}{|\rho|^2}-\frac{\gamma^2-1/4}{|\rho|^4}\bigr)
=\frac{2}{|\zeta'|^2}\cdot\frac{1/2}{|\rho|^4}$.
\end{proof}

\begin{remark}
\label{rem:cherm_not_identical}
The pointwise ratio is $(\gamma^2-\tfrac14)/(\gamma^2+\tfrac14)$,
approximately $0.9975005$ at the first zero: the two constants are close
but not identical. The identification of a recorded logarithmic
estimator $\hat S_2$ with $\beta_M$ requires RH, simplicity, convergence
and a limit, and is not asserted here.
\end{remark}

\subsection{The form $\mathcal A_0'$ is a L\'evy area}
\label{subsec:levy_area}

\begin{definition}
\label{def:A0}
With $M_\Omega(x)=\sum_{n\le x}\mu(n)\Omega(n)$ and
$L(x)=\sum_{n\le x}\mu(n)\log n$, the form recorded in separate work is
\begin{equation}
\mathcal A_0'(N)=\frac12\sum_{m<n\le N}\mu(m)\mu(n)
\bigl(\Omega(m)\log n-\log m\,\Omega(n)\bigr).
\label{eq:A0}
\end{equation}
\end{definition}

\begin{proposition}
\label{prop:levy_area}
\begin{equation}
\Delta\mathcal A_0'(N)=\frac{\mu(N)}{2}
\bigl[M_\Omega(N-1)\log N-\Omega(N)L(N-1)\bigr],
\label{eq:A0_increment}
\end{equation}
and consequently
\begin{equation}
\mathcal A_0'(N)=\mathfrak A_N(M_\Omega,L),
\qquad
\mathfrak A_N(X,Y)=\frac12\sum_{n\le N}\bigl(X(n-1)\,\Delta Y(n)-Y(n-1)\,\Delta X(n)\bigr),
\label{eq:levy_area}
\end{equation}
the discrete signed area swept by the planar path
$n\mapsto(M_\Omega(n),L(n))$.
\end{proposition}

\begin{proof}
The terms of \eqref{eq:A0} with $n=N$ give \eqref{eq:A0_increment}.
Summing, with $\Delta M_\Omega(n)=\mu(n)\Omega(n)$ and
$\Delta L(n)=\mu(n)\log n$, gives \eqref{eq:levy_area}.
\end{proof}

\begin{remark}
\label{rem:levy_area_meaning}
A L\'evy area is the antisymmetric second level of an iterated sum, so
it is neither a product of two summatory functions, which is the
symmetric part, nor a fixed-shift correlation. This is why the
dichotomy proposed in session discussion (either a product controlled by
RH, or a Chowla-type correlation) is false. Separate work records
$\Delta\mathcal A_0'(N)=\mu(N)N/(2\log N)+o(N/\log N)$.
\end{remark}

\begin{proposition}
\label{prop:F_of_s}
For $\Re s>2$ let
$F(s)=\sum_{n\ge2}\mu(n)n^{1-s}/\log n$. Then
\begin{equation}
F'(s)=1-\frac{1}{\zeta(s-1)},
\qquad
F(s)=\int_s^\infty\Bigl(\frac{1}{\zeta(u-1)}-1\Bigr)du,
\label{eq:F_of_s}
\end{equation}
so $F$ has logarithmic singularities at $s=1+\rho$.
\end{proposition}

\begin{proof}
Differentiate termwise:
$F'(s)=-\sum_{n\ge2}\mu(n)n^{1-s}=-(1/\zeta(s-1)-1)$. Since $F(s)\to0$ as
$s\to+\infty$, integrate from $s$ to $\infty$. Simple poles of
$1/\zeta(s-1)$ at $s-1=\rho$ integrate to logarithms.
\end{proof}

\begin{remark}
\label{rem:perron_mellin_obstacle}
The precise obstacle on the Perron--Mellin route is then to show that the
exact remainder of the Dirichlet series of $\mathcal A_0'$ continues far
enough and does not cancel these singularities.
\end{remark}

\subsection{Energies as aggregates of shifted correlations}
\label{subsec:energy_chowla}

\begin{lemma}
\label{lem:log_energy}
With $W_N(r)=\sum_{n=r}^{N}n^{-2}$,
\begin{equation}
\sum_{n\le N}\frac{M(n)^2}{n^2}
=\sum_{a\le N}\mu(a)^2W_N(a)
+2\sum_{h=1}^{N-1}\sum_{a\le N-h}\mu(a)\mu(a+h)\,W_N(a+h).
\label{eq:log_energy}
\end{equation}
\end{lemma}

\begin{proof}
Expand $M(n)^2$; the pair $(a,b)$ contributes $\sum_{n=\max(a,b)}^Nn^{-2}$.
\end{proof}

\begin{remark}
\label{rem:two_kernels}
Lemma~\ref{lem:energy_as_two_point} and Lemma~\ref{lem:log_energy} show
that both energies are weighted aggregates of shifted correlations. The
logarithmically averaged two-point theorem \cite{Tao2016} controls fixed
$h$; these aggregates need uniformity up to $h\asymp N$, and that gap is
not closed by naming the logarithmic theorem. The two recorded ratios
belong to different kernels, about $0.137$ for the area energy and about
$0.052$ for the truncated logarithmic energy, and are not comparable.
\end{remark}

\begin{remark}[Corrections ledger]
\label{rem:corrections_ledger}
Corrected in Version 5, in accordance with the archival principle:
the pushed form \eqref{eq:pushed_form}, previously written without the
incremental factors; the orientation basis \eqref{eq:canonical_F},
previously mixing reflections relative to $\mu$ with canonical signs; the
conditioning on $M(N)$ in Proposition~\ref{prop:reflection_fibre}; the
parity statement of Proposition~\ref{prop:pushed_form_odd}; the start of
the proof of Proposition~\ref{prop:log_exceptional}; the density claim of
Remark~\ref{rem:negative_result_density_of_states}; the scope of
Warning~\ref{warn:dichotomy}; and the crossover conjecture, withdrawn in
Remark~\ref{conj:twisted_orientation_bits}.
\end{remark}

\section{Strategy and the target theorem}
\label{sec:strategy}

\begin{remark}[Ordered strategy]
\label{rem:ordered_strategy}
\begin{enumerate}[label=(\roman*)]
\item Difference, not commutator: build $(U_h-I)D$, split off the
columns $d\mid h$ exactly by Lemma~\ref{lem:boundary}, and send the rest
to adaptive Haar.
\item Approximate observability: work in $V/\ker D_\Lambda$, measure the
optimal residual by pseudoinverse and the dual
\eqref{eq:operator_norm_dual}; do not attempt global identification,
which Thm.~\ref{thm:sublinear_nogo} excludes.
\item Irregular tree: choose cuts by changes of divisor density, levels
of $M$, centres and residue boundaries; exploit
\eqref{eq:cocycle_additivity}.
\item Local $p$-adic structure: couple
Prop.~\ref{prop:local_tables} with the factors $L_p$ of \ref{item:S4},
treating $p\nmid h$ as a universal factor and $p\mid h$ as diagonal
defects.
\item Excursions last: use them only for the reflection-odd component
and as one family of stopping blocks, subject to \ref{item:Q1}.
\end{enumerate}
\end{remark}

\begin{objective}[Adaptive boundary observability]
\label{obj:adaptive_boundary}
For fixed $h$, construct trees $\mathcal T_N$, multiscale subspaces
$W_N\subset H_{\mathrm{sf}}$ and rows $\Lambda_N$ adapted to CRT
boundaries, not to the signs of $\mu$, such that
\begin{equation}
\bigl\|P_{W_N}\bigl(A_h-L^{\mathsf T}B_{h,N}L\bigr)P_{W_N}\bigr\|_{\mathrm{op}}
=o(1),
\label{eq:obj_first}
\end{equation}
with $B_{h,N}$ the variational minimiser, and simultaneously
\begin{equation}
\bigl\|P_{W_N^{\perp}}\mu\bigr\|^2=o(N).
\label{eq:obj_second}
\end{equation}
Estimate \eqref{eq:obj_first} is to be proved by reducing all volumetric
contributions to zero through Theorem~\ref{thm:haar_crt} and controlling
only the boundary cocycles; \eqref{eq:obj_second} must come from prefix,
level and squarefree identities, not from probabilistic regularity.
Together they would give $C_N(h)=o(N)$.
\end{objective}

\begin{objective}[Variant with central divisor rows]
\label{obj:central_rows}
With
$F_{N,y}=\begin{bmatrix}V_y\\ C_N^{\mathrm{div}}\end{bmatrix}$,
truncated divisor equations together with centre rows fixed
independently of the orientations, prove
\begin{equation}
J_h^{\mathrm{off}}=F_{N,y}^{\mathsf T}B_{h,N}F_{N,y}+R_{h,N}
\label{eq:obj_central}
\end{equation}
with $B_{h,N}(d,e)=0$ unless $(d,e)\mid h$, prime-by-prime CRT
structure, $\|B_{h,N}\|=N^{o(1)}$, and
$\sup_{s\in\{\pm1\}^Z}|s^{\mathsf T}R_{h,N}s|=o(N)$, no definition using
$s_\mu$ or $C_N(h)$. Since $F_{N,y}s_\mu=0$ by
Theorem~\ref{thm:mobius_line}, this gives
$s_\mu^{\mathsf T}J_h^{\mathrm{off}}s_\mu=o(N)$; combined with the
independent intrablock objective
$\sum_j\sum_{n,n+h\in E_j}\mu(n)\mu(n+h)=o(N)$ it gives $C_N(h)=o(N)$.
\end{objective}

\subsection{Energy capture and the squeeze}
\label{subsec:squeeze}

\begin{proposition}[Exact energy of a leaf projection]
\label{prop:leaf_energy}
Let $W\subset H_{\mathrm{sf}}$ be the space of vectors constant on the
leaves $I$ of a partition of $\{1,\dots,m\}$ into blocks of consecutive
indices, with $k_I=|I|$. Then
\begin{equation}
\bigl\|P_{W^{\perp}}\mu\bigr\|^2
=\sum_I\Bigl(k_I-1-\frac{2}{k_I}\sum_{\substack{i<i'\\ i,i'\in I}}
\mu(a_i)\mu(a_{i'})\Bigr).
\label{eq:leaf_energy}
\end{equation}
Each summand is nonnegative and vanishes if and only if $\mu$ is
constant on the leaf.
\end{proposition}

\begin{proof}
$P_W\mu$ is the leaf average, so, using $\mu(a_i)^2=1$,
$\|P_W\mu\|^2=\sum_I k_I^{-1}\bigl(\sum_{i\in I}\mu(a_i)\bigr)^2
=\sum_I\bigl(1+\tfrac{2}{k_I}\sum_{i<i'}\mu(a_i)\mu(a_{i'})\bigr)$;
subtract from $\|\mu\|^2=\sum_Ik_I$. Nonnegativity and the equality case
are Cauchy--Schwarz on each leaf.
\end{proof}

\begin{remark}
\label{rem:haar_is_leaf}
The span of the constant and the Haar functions of a truncated tree is
exactly the space of vectors constant on the leaves of the truncation,
so Proposition~\ref{prop:leaf_energy} covers every multiscale subspace
admitted in Objective~\ref{obj:adaptive_boundary}.
\end{remark}

\begin{corollary}[The squeeze]
\label{cor:squeeze}
Condition \eqref{eq:obj_second} holds if and only if
\begin{equation}
\sum_I(k_I-1)-\sum_I\frac{2}{k_I}\sum_{\substack{i<i'\\ i,i'\in I}}
\mu(a_i)\mu(a_{i'})=o(m).
\label{eq:squeeze}
\end{equation}
Consequently:
\begin{enumerate}[label=(\roman*)]
\item singleton leaves contribute nothing, so \eqref{eq:obj_second}
holds trivially when all but $o(m)$ indices lie in singleton leaves; but
then $W_N$ is essentially all of $H_{\mathrm{sf}}$ and
\eqref{eq:obj_first} is the full observability problem, subject to
Theorem~\ref{thm:sublinear_nogo};
\item on leaves of size $k\ge2$ the condition requires the within-leaf
correlations of $\mu$ along consecutive squarefree integers to be
maximal, $\sum_{i<i'}\mu(a_i)\mu(a_{i'})=\binom{k}{2}$ up to $o$ of the
leaf mass: the opposite of the cancellation the objective is designed to
establish.
\end{enumerate}
\end{corollary}

\begin{proof}
Immediate from \eqref{eq:leaf_energy}.
\end{proof}

\begin{remark}[Long leaves are excluded unconditionally]
\label{rem:long_leaves}
For a leaf with $k_I\to\infty$, $\sum_{i\in I}\mu(a_i)$ is the sum of
$\mu$ over an integer interval of comparable length, since non-squarefree
integers contribute zero. For partitions to which the short-interval
estimate of \cite{MatomakiRadziwill2016} applies, in particular on
average over the offset of a grid, these sums are $o(k_I)$ in mean
square, hence $\|P_W\mu\|^2=o(m)$ and \eqref{eq:obj_second} fails by
the full amount. The input is used here only to refute.
\end{remark}

\begin{warning}
\label{warn:adaptive_objective_superseded}
Objective~\ref{obj:adaptive_boundary} is superseded as stated. Its
capture condition cannot be met by a subspace built independently of the
signs of $\mu$ except by taking essentially all of $H_{\mathrm{sf}}$,
and on leaves of size at least two it demands correlations of the wrong
sign. In addition, the claimed implication silently requires a norm
condition: writing $\mu=\mu_W+\mu_\perp$,
\begin{equation}
\mu_W^{\mathsf T}L^{\mathsf T}BL\,\mu_W=(L\mu_W)^{\mathsf T}B(L\mu_W),
\qquad L\mu_W=L\mu-L\mu_\perp,\qquad L\mu=D_\Lambda\mu\in\{0,e_1\},
\label{eq:hidden_norm}
\end{equation}
because $(D\mu)(r)=\sum_{d\mid r}\mu(d)=\mathbf 1[r=1]$; this term is
$o(N)$ only if $\|B_{h,N}\|\,\|L\|^2\,\|\mu_\perp\|^2=o(N)$. The Haar--CRT
vanishing of Theorem~\ref{thm:haar_crt} is unaffected, but it must be
used without an $L^2$ capture condition.
\end{warning}

\begin{remark}[What survives]
\label{rem:what_survives}
The engine common to both objectives is the annihilation
$D\mu=e_1$ of \eqref{eq:hidden_norm}. Objective~\ref{obj:adaptive_boundary}
wrapped it in a capture condition that is self-defeating by
Corollary~\ref{cor:squeeze}. Objective~\ref{obj:central_rows} uses it
directly, through $F_{N,y}s_\mu=0$ in Theorem~\ref{thm:mobius_line}, and
requires uniformity over the fibre rather than capture. It is therefore
the live target, and the boundary-cocycle estimate of
Theorem~\ref{thm:haar_crt} should be recast as a tool for bounding the
remainder $R_{h,N}$ there. Uniformity over all of $V$ is excluded by
Theorem~\ref{thm:sublinear_nogo}; uniformity over the fibre
$\{\pm1\}^Z$ is not.
\end{remark}

\begin{warning}
\label{warn:new_estimate_required}
The ordinary correlation $C_N(h)=o(N)$ at fixed $h\ne0$ is open; the
logarithmically averaged version is known \cite{Tao2016}, building on
\cite{MatomakiRadziwill2016}. Any inequality that yields $o(N)$
directly must therefore contain a genuinely new estimate, not a
conjugation or a rewriting. Remark~\ref{rem:conjugation_nogo} and
Warning~\ref{warn:path_reading_is_tautological} record two ways of
failing this test.
\end{warning}

\section{Negative results and failed routes}
\label{sec:negative_results}

\begin{remark}
\label{rem:negative_result_energy_scalar}
The scalar quadratic energy at a single level carries no local
information; the defect is the collapse to one number, not the weight.
\end{remark}

\begin{remark}
\label{rem:negative_result_height_weighting}
Statistics built from $M$ alone weight zeros by decaying factors in the
ordinate and therefore mix heights; no operation on $\log N$ repairs
this, since $\log N$ is conjugate to differences of ordinates, not to
ordinates.
\end{remark}

\begin{remark}
\label{rem:negative_result_jacobi}
The excursion decomposition passes to continued fractions and hence to a
Jacobi matrix \cite{Flajolet1980}, which is real symmetric and lands in
the orthogonal class, as in Rem.~\ref{rem:symmetry_class_obstruction}.
\end{remark}

\begin{remark}
\label{rem:negative_result_topology}
A cell-complex description of $M$ exists \cite{Bjorner2011} but does not
reach two-point cancellation: the complex is multiplicative while
shifted correlations are additive, and no map relates the complex at $n$
to that at $n+h$.
\end{remark}

\begin{remark}
\label{rem:negative_result_density_of_states}
Any operator whose spectrum is the set of ordinates must reproduce
$N(T)\sim(T/2\pi)\log(T/2\pi)$; the bounds \ref{item:S1} place
$\lambda_{\max}(D_N)$ at scale $N^{1/2+o(1)}$. This alone does not imply
a semicircular density, and the earlier claim to that effect is
withdrawn. What it does imply is that the first $N$ eigenvalues of $D_N$
cannot be $\gamma_1,\dots,\gamma_N$, since
$\gamma_N\sim2\pi N/\log N$ exceeds $C_\varepsilon N^{1/2+\varepsilon}$.
\end{remark}

\begin{remark}
\label{rem:negative_result_markov_tree}
A dyadic tree of midpoints of $|M|$ satisfies Chapman--Kolmogorov, hence
is Markov in the scale variable; a constraint at a coarse level
therefore gives no information at finer levels beyond the endpoint
values, and the degrees of freedom double at each refinement. The
product structure that makes the tree null computable is exactly the
property that forbids constraint propagation. This closes the route of
propagating a short-interval input down the tree. The same applies to
the Euler factorisation in the multiplicative direction: both natural
indexings factorise, so neither transports a constraint across scales.
\end{remark}

\begin{remark}
\label{rem:negative_result_omega_lower_bound}
Any hierarchy of bounds on the M\"obius walk must respect
$M(x)=\Omega_\pm(\sqrt x)$ with explicit constants
\cite{OdlyzkoteRiele1985}; no $\varepsilon$ in such a hierarchy is
removable.
\end{remark}

\begin{remark}
\label{rem:negative_result_halves_crt}
Splitting an interval into two halves keeps an $O(1)$ CRT error in each
half; it does not reduce it. A left--right difference can cancel the main
term when the lengths agree, but leaves a difference of boundary errors
of size $O(1)$ with no sign and no uniform cancellation. For halves
defined in squarefree time the set is no longer an arithmetic interval
and the mask must be inserted, which makes the CRT factorisation less
clean, not more.
\end{remark}

\begin{remark}[The irregular route, closed as a family]
\label{rem:irregular_family_closed}
Every divisor-to-shift transfer that uses only $o(N)$ rows of $D$
(Corollary~\ref{cor:sublinear_quantitative}), compression by prime
valuations (Proposition~\ref{prop:valuation_nogo}), irregular Haar
without an additional arithmetic estimate of the boundary cocycles
(Proposition~\ref{prop:tree_independent}), or separable approximations of
rank $o(N)$ (Remark~\ref{rem:blind_class_meaning}) has macroscopic
operator error on the squarefree subspace and cannot imply
$C_N(h)=o(N)$. The quantitative progress available remains for
logarithmic averages \cite{Tao2016, Pilatte2023}.
\end{remark}

\begin{heuristic}
\label{heur:collective_crt}
The one door left open inside this family is a collective signed
cancellation of the CRT boundary error,
\begin{equation}
\sum_{d,e}c_dc_e\,\epsilon_{d,e,h},
\label{eq:collective_crt}
\end{equation}
for coefficients tied specifically to $\mu$. Obtaining $o(N)$ for this
form already requires control of the global parity of large primes,
which is exactly the content that local factors, central rows and the
squarefree mask do not contain.
\end{heuristic}

\section{Open questions}
\label{sec:open_questions}

\begin{enumerate}[label=(Q\arabic*)]
\item \label{item:Q1}
Does $\widehat G$ inherit the tensor factorisation
\eqref{eq:tensor_factorisation_standing} from $G$? The excursion blocks
are defined by the zero set of $M$ and respect no multiplicative
structure, so clean inheritance is not expected; partial inheritance
restricted to the local factors $L_p$ with small $p$ would be the
interesting outcome. Everything in
Sections~\ref{sec:scale_resolved}--\ref{sec:strategy} depends on it.
\item Prove a norm estimate for the boundary CRT cocycle operator of
Theorem~\ref{thm:haar_crt}, restricted to adaptive Haar subspaces and
after extracting the exact local $p$-adic modes of
Prop.~\ref{prop:local_tables}, respecting
Warning~\ref{warn:parseval_nogo}. Following
Warning~\ref{warn:adaptive_objective_superseded}, this estimate is now
sought as a bound for $R_{h,N}$ in Objective~\ref{obj:central_rows},
not as the first step of the superseded
Objective~\ref{obj:adaptive_boundary}.
\item Sharpen Prop.~\ref{prop:log_exceptional}: is the exceptional set
actually bounded, or can $\log_2N$ be attained?
\item Test the prediction that $1/p(1)$ is the minimum of the
divisor-resolved sine-kernel fraction.
\item \label{item:Q5}
Determine the growth of $Z(N)$. Lemma~\ref{lem:left_right} and every
statement with an $O(hZ)$ error depend on $Z=o(N)$, which is not proved
here.
\item Close the non-attained-supremum case of
Rem.~\ref{rem:extremal_abscissa}.
\item Decide whether the non-commutative deformation of
Warning~\ref{warn:dichotomy} escapes \ref{item:S8} in the scaling window
$\vartheta\asymp\vartheta_c(N)$.
\end{enumerate}

\section*{Notebook History}
\label{sec:notebook_history}

Version 1:
Three-constraint model and exact bridge second moment; reflection fibre
and the canonical parameter-free null; divisor-defect Gram form pushed
to orientation bits; blindness theorem; scale-resolved statistic with
exact mean and variance; negative results.

Version 2:
Abstract added; introduction restructured into what is proved, what is
input, and the methodological point; standing results listed with status;
sine-kernel section rewritten around the non-tautology criterion, with
three recoveries of the sine-kernel coefficient from measured
quantities; extremal-abscissa gap narrowed.

Version 3:
Operator layer added. Three spaces and their information ledger;
divisor--shift boundary lemma with interval mass bound $1$; exact CRT
pullback; algebraic Haar on arbitrary trees with exact vanishing of the
volumetric CRT term and parent--children cocycle additivity; exact
observability criterion with optimal constant, operator-norm dual and
semidefinite formulation; no-go for sublinear divisor rows; extension of
the blindness theorem to half-constant kernels, with the shift shown to
escape it; central dichotomy lemma; triangular central divisor matrix
with the logarithmic exceptional set, isolating the M\"obius line;
intrablock second-difference and aggregate identities, discrete coarea,
and a counterexample showing the centre does not determine a lag; exact
local tables at each prime; collapse theorem for multiplicative diagonal
weights and the resulting dichotomy; target objectives stated.

Version 4:
Audit of the target objective. Conditioning of the two-point expansion
on central data; redundancy of quarter and octave evaluations; exact
leaf-energy identity for leaf projections in the squarefree space; the
squeeze corollary showing that the capture condition of the adaptive
boundary objective forces essentially the whole squarefree space or
maximal positive within-leaf correlation; exclusion of long leaves by
short-interval cancellation; the hidden norm condition in the claimed
implication. The adaptive boundary objective is marked superseded; the
central-row objective, built on $D\mu=e_1$ and uniformity over the
fibre, is the live target.

Version 5:
External audit incorporated after line-by-line verification. Pushed form
corrected to $F^{\mathsf T}GF$; single orientation basis; conditioning on
$M(N)$ made explicit; single-reflection parity corrected; exceptional-set
proof started at $z_1=2$; semicircle claim and general weight dichotomy
withdrawn; crossover conjecture withdrawn by the gauge obstruction, with a
flux question and a crossover-scale heuristic in its place. New: the
energy kernel as Montgomery's weight at $\alpha=0$ with explicit density
and tail; exact fourth-order relation between the Hermitian collision
constant and the M\"obius energy constant; $\mathcal A_0'$ as a discrete
L\'evy area, with its increment and the function $F(s)$; logarithmic
energy as an aggregate of shifted correlations; compressed selection to
at most $2N$ orientations and unique selection after
$\lfloor\log_2N\rfloor+1$ further rows.

Version 6:
Uncompressed irregular route executed. Boundary lemma restricted to the
interior operator, with the boundary rows made explicit; maximal rank of
the interior operator; exact $1\to1$, $\infty\to\infty$ and
Hilbert--Schmidt norms; residue-boundary form of each column with the
optimal variation bound; tree-independent Haar energy; quantitative
observability no-go with constant $\mathfrak S_h/4$; signed local sums
and the Mirsky expansion; failure of the shift to descend to valuations;
maximal blind class as the additively separable kernels, with the
distance of the shift and the Eckart--Young bound; global and block
second-difference identities with explicit boundary terms. The irregular
route is closed as a family; the collective CRT cancellation is recorded
as the remaining heuristic.

\begin{thebibliography}{99}

\bibitem{AtasEtAl2013}
Y.~Y.~Atas, E.~Bogomolny, O.~Giraud and G.~Roux,
\emph{Distribution of the ratio of consecutive level spacings in random
matrix ensembles},
Physical Review Letters \textbf{110} (2013), 084101.

\bibitem{Bjorner2011}
A.~Bj\"orner,
\emph{A cell complex in number theory},
arXiv:1101.5704.

\bibitem{Flajolet1980}
P.~Flajolet,
\emph{Combinatorial aspects of continued fractions},
Discrete Mathematics \textbf{32} (1980), 125--161.

\bibitem{Gonek1989}
S.~M.~Gonek,
\emph{On negative moments of the Riemann zeta-function},
Mathematika \textbf{36} (1989), 71--88.

\bibitem{Johansson1998}
K.~Johansson,
\emph{On fluctuations of eigenvalues of random Hermitian matrices},
Duke Mathematical Journal \textbf{91} (1998), 151--204.

\bibitem{Lagarias2002}
J.~C.~Lagarias,
\emph{An elementary problem equivalent to the Riemann hypothesis},
American Mathematical Monthly \textbf{109} (2002), 534--543.

\bibitem{MatomakiRadziwill2016}
K.~Matom\"aki and M.~Radziwi\l{}\l{},
\emph{Multiplicative functions in short intervals},
Annals of Mathematics \textbf{183} (2016), 1015--1056.

\bibitem{Mehta2004}
M.~L.~Mehta,
\emph{Random Matrices}, third edition,
Elsevier, Amsterdam, 2004.

\bibitem{Mirsky1949}
L.~Mirsky,
\emph{Arithmetical pattern problems relating to divisibility by $r$th
powers},
Proceedings of the London Mathematical Society \textbf{50} (1949),
497--508.

\bibitem{Montgomery1973}
H.~L.~Montgomery,
\emph{The pair correlation of zeros of the zeta function},
in: Analytic Number Theory,
Proceedings of Symposia in Pure Mathematics \textbf{24},
American Mathematical Society, 1973, 181--193.

\bibitem{Odlyzko1987}
A.~M.~Odlyzko,
\emph{On the distribution of spacings between zeros of the zeta
function},
Mathematics of Computation \textbf{48} (1987), 273--308.

\bibitem{OdlyzkoteRiele1985}
A.~M.~Odlyzko and H.~J.~J.~te Riele,
\emph{Disproof of the Mertens conjecture},
Journal f\"ur die reine und angewandte Mathematik \textbf{357} (1985),
138--160.

\bibitem{PandeyMehta1983}
A.~Pandey and M.~L.~Mehta,
\emph{Gaussian ensembles of random Hermitian matrices intermediate
between orthogonal and unitary ones},
Communications in Mathematical Physics \textbf{87} (1983), 449--468.

\bibitem{Pilatte2023}
C.~Pilatte,
\emph{Improved bounds for the two-point logarithmic Chowla conjecture},
arXiv preprint, 2023.

\bibitem{Robin1984}
G.~Robin,
\emph{Grandes valeurs de la fonction somme des diviseurs et hypoth\`ese
de Riemann},
Journal de Math\'ematiques Pures et Appliqu\'ees \textbf{63} (1984),
187--213.

\bibitem{Tao2016}
T.~Tao,
\emph{The logarithmically averaged Chowla and Elliott conjectures for
two-point correlations},
Forum of Mathematics, Pi \textbf{4} (2016), e8.

\bibitem{Titchmarsh1986}
E.~C.~Titchmarsh,
\emph{The Theory of the Riemann Zeta-Function}, second edition,
revised by D.~R.~Heath-Brown,
Oxford University Press, 1986.

\end{thebibliography}

\label{sec:bibliography}

\end{document}
